%%%%%%%%%%%%%%%%%%%%%%%%%%%%%%%%%%%%%%%%%%%%%%%%%%%%%%%%%%%%%%%%%%%%%%%%%%%%%%%%
%%                           ONDER DE LOEP ARTIKEL                            %%
%%%%%%%%%%%%%%%%%%%%%%%%%%%%%%%%%%%%%%%%%%%%%%%%%%%%%%%%%%%%%%%%%%%%%%%%%%%%%%%%

\artikelfoto{Statistiek in de derde graad humane wetenschappen}{}{Jan De Neve \\ Tom Loeys \\ Dennis Presotto  \\ Luc Van den Broeck}{img/Lkopfoto.jpg}

\startcontents[inhoudloep]	% om inhoud voor het loep artikel te beginnen

\begin{multicols}{2}
	
	\inhoudstafelloep			% om inhoudstafel voor het loep artikel te tonen
	
	%% Gebruik \section{titel} en \subsection{titel} om genummerde tussentitels te maken
	\section{Inleiding}
        \subsection*{Een woordje vooraf}
        Deze loep is het werk van vier auteurs. Vanuit de redactie werkten Luc Van den Broeck en Dennis Presotto aan dit artikel. Luc is hierin een gevestigde waarde. Dennis is een nieuw lid sinds 2022. Hij is leerkracht wiskunde en statistiek aan de Stedelijke Humaniora Dilsen. Daarnaast werkten ook Jan De Neve en Tom Loeys mee aan dit artikel als gastauteurs. Beiden zijn statistici en doceren verschillende opleidingsonderdelen statistiek aan de UGent in de faculteit Psychologie en Pedagogische Wetenschappen.

        Uitwiskeling 34/1 nam eerder al \emph{Statistiek buiten het boekje} onder de loep. Verschillende idee\"en uit dat artikel zullen hier terugkomen. Toch zijn er ook verschillen. Zo lag destijds de focus op uitbreidingen in de diepte (centrale limietstelling) en in de breedte (andere verdelingen en formules). Het specifieke doelpubliek van humane wetenschappers en de evolutie naar de GRexit (zoals besproken in Uitwiskeling 38/3) zorgen voor een verhoogd belang van context en begripsvorming.
        
	\subsection*{Waarom extra statistiek?}
        Hoewel op het moment van schrijven de eindtermen voor de derde graad nog niet definitief zijn, kunnen we moeilijk ontkennen dat er enkele veranderingen aankomen. 
        
        Deze loep focust op het nieuwe pakket eindtermen in de statistiek dat vanaf het schooljaar 2023-2024 verplicht zal worden in de richtingen humane wetenschappen, welzijnswetenschappen en Rudolf Steiner pedagogie. Alvorens in te gaan op de eindtermen die hier (vermoedelijk) opgelegd zullen worden, is het misschien leuk om te weten waarom die er eigenlijk komen en waarom dit specifiek bij deze richtingen gebeurt. Daarbij zullen we iets meer focussen op de studierichting humane wetenschappen. Enerzijds zitten in deze richting beduidend meer leerlingen dan in de andere twee richtingen, anderzijds hebben de schrijvers meer ervaring in deze richting dan in de twee andere. In één zin zouden we kunnen stellen dat de veranderingen in deze richtingen gebeuren op vraag van het hoger onderwijs. Het antwoord in iets meer zinnen volgt nu. 
        
      
        Statistiekvakken komen in verschillende opleidingen in het hoger onderwijs voor. In het bijzonder focussen we hier op de opleidingen psychologie, toegepaste psychologie, criminologische wetenschappen, pedagogische wetenschappen, sociologie, politieke wetenschappen en communicatiewetenschappen. Deze verschillende opleidingen hadden in het academiejaar 2021-2022 (volgens \url{onderwijskiezer.be}) een gelijkaardige instroom uit het secundair onderwijs: de grootste groep instromende studenten kwam voor elk van deze opleidingen uit de humane wetenschappen. Dit is het meest uitgesproken bij de academische bachelor-opleiding psychologie, waarbij 2394 ingeschreven studenten uit de richting humane wetenschappen kwamen. Dat was meer dan een derde van alle inschrijvingen van eerste bachelor psychologie en ongeveer 9\% van alle leerlingen die in het schooljaar 2020-2021 afstudeerden als humane wetenschapper. 
        
        Qua studierendement scoren de humane wetenschappers echter slechter dan verschillende andere instromers. Vaak hebben die andere studenten meer uren wiskunde gehad in het middelbaar onderwijs. Dit toont uiteraard geen oorzakelijk verband aan, maar het is op zijn minst opvallend.

        Wanneer we het hebben over studierendement, dan moeten we het ook hebben over de `knip'. In het regeerakkoord van de Vlaamse regering van 30 september 2019 werd er gesproken over een `sterkere knip tussen bachelor en master' waarbij het behalen van bepaalde bachelorvakken binnen een vastgelegde tijd een voorwaarde werd voor inschrijving in een masteropleiding. Dit voornemen kende wat tegenstand en werd hervormd naar een `knip' na het tweede bachelorjaar zoals goedgekeurd in het decreet van 13 juli 2022. Concreet zal vanaf het academiejaar 2023-2024 elke student op het einde van het tweede bachelorjaar geslaagd moeten zijn op alle vakken van het eerste bachelorjaar. De regering neemt hierbij de visie aan dat studenten de moeilijkere basisvakken (zoals statistiek) uit het eerste bachelorjaar niet mogen laten liggen wanneer die relevant zijn voor het vervolg van hun opleiding (zoals bij psychologie). De universiteiten en hogescholen kunnen voor individuele studenten van deze regel afwijken, maar dit is enkel toegelaten in uitzonderlijke gevallen. 

        Uit verschillende studies blijkt trouwens dat angst voor wiskunde, en in het bijzonder voor statistiek, leeft bij een groot deel van de studerende bevolking, zie bijvoorbeeld (Siew et al., 2019). Dat beeld wordt tevens bevestigd door het grote aanbod aan personen en kantoren die bijlessen statistiek aanbieden. De nieuwe eindtermen zouden ervoor moeten zorgen dat instromende studenten een bredere basiskennis statistiek hebben, hetgeen hopelijk de angst ervoor kan onderdrukken. 
        
        Een laatste argument dat soms opgeworpen wordt voor het invoeren van de eindtermen statistiek is het opkrikken van het imago van de richting humane wetenschappen. Voor sommige leerlingen was in het verleden de keuze voor deze studierichting eerder een negatieve keuze om wiskunde of wetenschappen te ontlopen dan een positieve keuze voor het humane en sociale aspect van de richting. Bij de groep waar de keuze voor deze studierichting wel positief was, treffen we regelmatig leerlingen aan die wiskundig intrinsiek minder sterk zijn dan in andere doorstroomrichtingen. Voor beide groepen zou deze extra statistiek een oplossing kunnen bieden. Dat lezen we bijvoorbeeld in het artikel uit De Standaard van 6 september 2021 (Maenhout, 2021).
        
        \subsection*{Welke eindtermen worden behandeld?}
        Als je al eens gekeken hebt naar de eindtermen statistiek, dan is het je waarschijnlijk opgevallen dat de eindtermen statistiek als aanvulling dienen op het basispakket wiskunde voor de doorstroomfinaliteit. Verschillende eindtermen die reeds bestaan in het vak wiskunde worden verdergezet, maar er zijn eveneens nieuwe onderwerpen aan toegevoegd. Het bestaan als aanvullingspakket benadrukt ook dat wiskundige of statistische eindtermen kunnen vari\"eren tussen de studierichtingen. Zo zou in het verleden een humane wetenschapper een basispakket van (meestal) 3 uur wiskunde per week volgen in de derde graad en dit pakket was hetzelfde voor leerlingen uit een richting economie-moderne talen. Met de vernieuwde eindtermen delen deze richtingen nog steeds eenzelfde basispakket wiskunde, maar ze worden aangevuld met respectievelijk specifieke eindtermen \emph{statistiek} en  \emph{uitgebreide wiskunde in functie van economie}. Deze \emph{uitgebreide wiskunde} kent ook varianten voor andere studiedomeinen. De algemene visie hierachter is dat we leerlingen beter willen voorbereiden op de wiskunde en statistiek die relevant is in typische vervolgstudies om zo de slaagkansen en interesse te verhogen. 
        
        Hieronder bespreken we in meer detail de eindtermen zoals terug te vinden op \url{onderwijsdoelen.be} op het moment van schrijven (en dus voor de mogelijke aanpassingen van 2023). De aandachtige lezer zal daarbij opmerken dat verschillende, maar niet alle eindtermen ook elders terugkomen. Zo zullen bijvoorbeeld ook economisten en wetenschappers in contact komen met het toetsen van hypotheses.        

        Binnen \emph{statistiek} vinden we de volgende vier eindtermen:
 \begin{kader}{Eindterm}{Eindtermen statistiek}

        
        \begin{itemize}
        
        \item 6.1.1 De leerlingen gebruiken combinaties en de binomiale verdeling in betekenisvolle situaties.
        \item 6.1.2 De leerlingen leggen in betekenisvolle situaties de betekenis van betrouwbaarheidsniveau en betrouwbaarheidsinterval uit.
        \item 6.1.3 De leerlingen toetsen hypotheses.
        \item 6.1.4 De leerlingen analyseren grote datasets met behulp van statistische software in functie van een statistisch onderzoek.
    \end{itemize}
\end{kader}
    In de volgende paragrafen vind je meer details over eindtermen 6.1.1, 6.1.2 en 6.1.3. We eindigen deze loep met een paragraaf over mogelijke vakoverschrijdende projecten. Zulke projecten kunnen een invulling geven aan eindterm 6.1.4, al zal het hier zelden over echt grote datasets gaan. Het werken met heel grote datasets via ICT lijkt ons echter relevant voor meerdere studierichtingen.  Wanneer er meer duidelijkheid is over de eindtermen en de inzetbare voorkennis ICT van de instromende leerlingen uit de tweede graad wordt dit misschien het onderwerp van een toekomstige loep.

    In het algemeen valt alvast op dat er een aandachtsverschuiving is. Waar vroeger voornamelijk gekeken werd naar de normale verdeling (eindtermen wiskunde ASO voor de modernisering: 33, 34, 35 en 36) en waar de binomiale verdeling enkel voor de richtingen met meer wiskunde werd behouden (specifieke eindtermen wiskunde ASO: 17), krijgen beide nu aandacht. De binomiale verdeling vinden we terug in eindterm 6.1.1 van \emph{statistiek}; de normale verdeling staat voor alle richtingen met doorstroomfinaliteit bij de wiskundige eindterm 6.13. 
    
    Ongeacht de keuze van de school om de eindtermen wiskunde en statistiek in twee afzonderlijke vakken of in een gezamenlijk vak na te streven, vallen deze eindtermen voor de normale en binomiale verdeling moeilijk van elkaar los te koppelen. Beide verdelingen zijn relevant voor het toetsen van hypotheses en het opstellen van betrouwbaarheidsintervallen. In de volgende paragrafen vind je over deze statistische technieken meer uitleg en voorbeelden. Hierbij trachten we deze concepten in te voeren met zo weinig mogelijk wiskundige formules en proberen we vooral het denkkader te duiden. 

\subsection*{Indeling van deze loep }

In paragraaf \ref{S_2} starten we met het toetsen van hypotheses. Om tot het toetsen te komen worden eerst op informele wijze de stappen van een binomiaaltoets doorlopen door middel van lesactiviteiten. Dit geeft de leerlingen een eerste beeld van \emph{wat} een hypothesetoets is en \emph{waarom} we deze nodig hebben. 

In paragraaf \ref{S_3} gaan we dieper in op het kansmodel dat centraal staat in de binomiaaltoets: de binomiale verdeling. De aanbeveling om kansrekening tot een minimum te beperken, impliceert niet dat alle kansrekening moet geschrapt worden. Het is belangrijk om de rol van kansmodellen binnen inductieve statistiek te illustreren. Door deze echter pas te introduceren nadat de leerlingen informeel hebben kennisgemaakt met de binomiaaltoets, zou het duidelijker moeten zijn waarom we deze kansverdeling nodig hebben. 

Vervolgens formaliseren we de binomiaaltoets in paragraaf \ref{S_4}. In paragraaf \ref{S_5} maken we de overstap naar betrouwbaarheidsintervallen. In paragraaf \ref{S_6} sluiten we met vakoverschrijdende projecten die de rol en de meerwaarde van statistiek binnen de empirische cyclus demonstreren. We overlopen daarbij ook het GAISE-rapport dat verschillende aanbevelingen geeft voor het onderwijzen van statistiek.

\section{Toetsen van hypotheses} \label{S_2}

Het toetsen van hypotheses is een voorbeeld van \emph{inductieve of verklarende statistiek}. Meer specifiek laat het toe om algemene besluiten te formuleren over een volledige populatie op basis van data in een steekproef. We starten deze paragraaf met een korte lesactiviteit uit (Rossman, 2008) waarbij leerlingen zonder het te beseffen beroep doen op de manier van redeneren die gebruikt wordt bij een hypothesetoets. 

\subsection*{Een inleidend experiment}
    
We starten met het opgooien van twee dobbelstenen. Het zijn bijzondere dobbelstenen want de ene dobbelsteen heeft vijf op elke zijde, terwijl de andere dobbelsteen bestaat uit enkel twee of zes (je kunt die zelf aanmaken met blanco dobbelstenen maar je kunt evengoed een kaartspel gebruiken voor deze lesactiviteit). Vervolgens voer je verschillende worpen uit en schrijf je de som van de ogen op het bord. De leerlingen weten niet dat het bijzondere dobbelstenen zijn, maar ze zullen dit na een aantal worpen snel achterhalen. Om de leerlingen extra te betrekken kun je ze eventueel belonen als de som van de ogen even is. 
    
De som van de ogen kan maar twee waarden aannemen: zeven of elf. Naarmate je meer worpen uitvoert, zal je de reactie van de leerlingen zien veranderen. Bij de eerste worpen is er niets aan de hand, maar na een aantal worpen zullen de leerlingen argwanend worden en aangeven dat de dobbelstenen niet eerlijk zijn. Bij hun argumentatie komt het principe van de hypothesetoets aan bod: als de dobbelstenen eerlijk zijn (als m.a.w. de zogenaamde \emph{nulhypothese} waar is), dan is het weinig waarschijnlijk dat we een lange reeks met enkel zeven en elf observeren. Omdat de leerlingen net een lange reeks van zeven en elf observeren, besluiten ze dat de dobbelstenen niet eerlijk zijn. 


Je zou kunnen stellen dat de leerlingen gebruik maken van een \emph{inductieve} redenering aangevuld met een \emph{modus tollens} argument ($P$ impliceert $Q$ en $Q$ is niet waar. Daarom is $P$ ook niet waar). Deze terminologie hoeft niet expliciet bij de leerlingen gebruikt te worden, maar deze logische regel helpt wel om het principe van hypothese testen beter te begrijpen. Als de nulhypothese (propositie $P$) waar is, verwacht je data in de steekproef (propositie $Q$) in lijn hiermee. De steekproef blijkt niet aan de verwachtingen te voldoen (propositie $Q$ is niet waar). Op basis daarvan verwerpen de leerlingen de nulhypothese (propositie $P$ is niet waar). Uiteraard ligt het bij statistiek nog iets complexer, omdat er steeds onzekerheid met onze uitspraken gepaard gaat. 

Gezien we op basis van data in de steekproef, een uitspraak doen over de volledige populatie, spreken we hier van \emph{inductie}. Inductief redeneren is vaak niet evident omdat we binnen de wiskunde meestal gebruik maken van deductie. 


%\subsection*{Een tweede experiment}

%Voorgaand voorbeeld illustreert hoe de manier van redeneren bij een hypothesetoets vrij natuurlijk tot stand komt: op basis van geobserveerde data in een steekproef wordt gekeken of er bewijskracht gevonden kan worden tegen de nulhypothese over een fenomeen op populatieniveau. Toch wordt dit door velen als moeilijk ervaren, zeker als de hypothesetoets in haar abstracte vorm wordt uitgelegd. Het voorzien van context kan dit probleem oplossen. 

%We illustreren dit met een tweede voorbeeld uit (Rossman, 2008) waar \emph{modus tollens} centraal staat. In de abstracte versie van de oefening krijgen de leerlingen kaarten te zien. Elke kaart heeft een letter op de ene zijde en een getal op de andere zijde. De leerlingen moeten nagaan of volgende regel geldt: de kaarten met een klinker hebben een even getal op de andere zijde. Vervolgens krijgen ze vier kaarten te zien: A, B, 6 en 7. Welke kaarten moeten de leerlingen omdraaien om na te gaan of de regel geschonden is? 

%De kans is groot dat leerlingen A en 6 aanduiden (Wason en Johnson-Laird, 1972), terwijl het A en 7 moet zijn. De keuze voor 7 volgt immers uit een gekende regel uit de logica, de \emph{contrapositie}-regel (die stelt dat de proposities $P \Rightarrow Q$ en $\neg Q \Rightarrow \neg P$ equivalent zijn). In dit experiment wordt dit: als er geen even getal (m.a.w. een oneven getal) op de ene kant staat dan staat er geen klinker (m.a.w. een medeklinker) op de andere kant.

%Wanneer we hetzelfde probleem voorzien van context, zullen leerlingen het makkelijker vinden om tot een juist antwoord te komen. De nieuwe variant van de oefening gaat over een feestje waarop de volgende alcoholregel geldt: als je minderjarig bent dan drink je niet alcoholisch. We controleren vier personen: een dertigjarige, een zestienjarige, iemand die bier drinkt en iemand die cola drinkt. Welk van deze vier controles is nu geschikt om te testen of de alcoholregels gerespecteerd zijn?

%Leerlingen zullen nu makkelijker tot het juiste antwoord komen: ze moeten nagaan welke leeftijd de bierdrinker heeft en wat de zestienjarige drinkt. Ook dit is een toepassing van de contrapositie-regel. De uitspraak \emph{uit minderjarig volgt niet alcoholisch} is immers equivalent met de uitspraak \emph{uit alcoholisch volgt meerderjarig}.

In de volgende paragraaf tonen we aan hoe hypothesetoetsen gelijkaardig gebruik maken van een modus tollens-argument. Hierbij moet ook nog eens onzekerheid in rekening gebracht worden. Daarom zijn hypothesetoetsen intuïtief niet eenvoudig. Door context te voorzien die de leerlingen begrijpen, wordt de opdracht eenvoudiger en kunnen ze zelf het principe naar voor schuiven (zoals in het voorbeeld met de dobbelstenen). We passen dit idee toe om het begrip \emph{binomiaaltoets} te introduceren. 

\subsection*{De binomiaaltoets}

In het voorbeeld hieronder staan eerst de context en de onderzoeksvraag centraal zodat we de leerlingen (en hun intu\"itie) kunnen betrekken. Aansluitend bij de studierichting humane wetenschappen baseren we ons op een psychologisch onderzoek van Hamlin et al. (2007), een wetenschappelijk artikel uit een reeks waarbij onderzoekers op empirische wijze het moraliteitsbesef bij baby's onderzoeken. 

\begin{kader}{}{Context en onderzoeksvraag}

De capaciteit om anderen in te kunnen schatten is essentieel om met elkaar te kunnen omgaan. We moeten in staat zijn om acties en intenties van anderen in te schatten om vervolgens een beslissing te nemen over wie goede of kwade bedoelingen heeft. Volwassenen zullen anderen snel taxeren en evalueren op basis van zowel gedrag als lichamelijke eigenschappen. De ontwikkeling van dit taxatievermogen is echter niet volledig duidelijk. Om hier meer inzicht in te krijgen, wordt een studie opgezet waarin de morele intuïtie van baby's wordt onderzocht. Centraal staat de onderzoeksvraag:
\vspace{.2cm}
\\
\emph{Verkiezen baby's een pop die een goede daad stelt boven een pop die een slechte daad stelt?}
\end{kader}
\end{multicols}
\afbeelding[17.5cm]{img/LHamlinFiguur1.png}{Links: tafereel waarbij de helper (de gele driehoek) de klimmer helpt. Rechts: tafereel waarbij de lastpost (het blauwe vierkant) de klimmer tegenwerkt. Figuur overgenomen uit (Hamlin et al., 2007).}{fig:Hamlin1} 
\begin{multicols}{2}
De onderzoeksvraag heeft betrekking op een grotere groep baby's, ook wel de populatie genoemd. In het artikel wordt niet expliciet weergegeven wat deze groep is. Omdat men fundamentele eigenschappen onderzoekt, zou je kunnen stellen dat de onderzoeksvraag betrekking heeft op alle baby's die bijvoorbeeld tien maanden oud zijn.

We geven in een kader het ontwerp van de studie mee omdat statistiek meer is dan alleen bewerkingen uitvoeren op data. Minstens even belangrijk is het begrijpen en kritisch beoordelen van de manier waarop data tot stand zijn gekomen. Door voldoende stil te staan bij de context, de onderzoeksvraag en het ontwerp van de studie, kun je de onderzoekende attitude bij de leerlingen stimuleren. 

Na het kader geven we in een lesactiviteit enkele vragen en antwoorden die een klasgesprek hierrond kunnen stimuleren.

\begin{kader}{}{Ontwerp van de studie}

Zestien baby's van tien maanden oud namen deel aan de studie. Zij vormen de steekproef uit de hierboven beschreven populatie. De baby's kregen een poppenspel te zien waarbij een houten speelgoedpop (de `klimmer') een berg wil opklimmen. Bij sommige taferelen werd de klimmer geholpen door een andere houten speelgoedpop (de `helper'), terwijl in andere taferelen een derde houten speelgoedpop (de `lastpost') de klimmer tegenwerkt. \autoref{fig:Hamlin1} illustreert deze taferelen. Nadat de baby's verschillende van deze taferelen hadden gezien, werd hen de mogelijkheid gegeven om te kiezen tussen de helper of de lastpost: beide speelgoedpopjes werden aangeboden en men keek naar welke pop de baby greep.
\end{kader} 

\end{multicols}

\afbeelding[10cm]{img/UWbabies.jpg}{}{}

\clearpage

\begin{lesactiviteit}{Klasgesprek: evaluatie van het ontwerp van de studie}
	
\begin{vraagenantwoord}
    \vraag{Als je kijkt naar de opstelling uit \autoref{fig:Hamlin1}, zijn er dan naast de rol van elke pop (de helper en de lastpost) andere elementen die de keuze van de baby kunnen beïnvloeden?  }
    \antwoord{De helper en de lastpost verschillen in kleur en vorm. Indien baby's bijvoorbeeld een voorkeur hebben voor geel of voor een driehoek, dan kan de helper hierdoor vaker of minder vaak gekozen worden, ook indien de baby's geen voorkeur (of notie) hebben voor de daad die de pop stelt. Om deze mogelijke effecten te vermijden, hebben de onderzoekers de rol van de poppen gebalanceerd. Voor de helft van de baby's was de gele driehoek de helper, terwijl voor de andere helft deze pop de lastpost was.}
    
    \vraag{De baby's zitten op de schoot bij een van de ouders wanneer ze kunnen kiezen tussen beide poppen. Zou dit een invloed kunnen hebben op de keuze? }
    \antwoord{De ouder kan onbewust de keuze be\"invloeden door zich bijvoorbeeld meer te richten naar een van beide poppen. Om dit effect te vermijden, moest de ouder zijn/haar ogen sluiten bij de keuze.}

    \vraag{Een onderzoeker bood beide poppen aan zodat de baby er een kon grijpen. Zou dit een invloed kunnen hebben op de keuze? }
    \antwoord{Net zoals de ouder kan de onderzoeker ook een onbewuste invloed uitoefenen op de keuze door een pop dichter aan te bieden. Om dit te vermijden wist de onderzoeker niet welke rol de poppen hebben gespeeld (en was het een andere onderzoeker die het poppenspel speelde).}
\end{vraagenantwoord}
	
\end{lesactiviteit}

\begin{multicols}{2}

Via de bovenstaande lesactiviteit wordt het duidelijk dat we bij een goed ontwerp van de studie zoveel mogelijk andere oorzaken die de keuze van pop beïnvloeden, trachten uit te sluiten. Op die manier kunnen we overtuigender stellen dat de keuze gelinkt is aan de daad die de pop heeft gesteld. 

\begin{kader}{}{De data}

Veertien van de zestien baby's kozen de helper. 
\end{kader}

Nu kunnen we de leerlingen terug aan het werk zetten. De onderzoeksvraag heeft betrekking op baby's in het algemeen (de populatie) terwijl we data hebben voor slechts zestien baby's (de steekproef). Zijn veertien keuzes voor de helper in de steekproef voldoende bewijs om te besluiten dat de helper in de populatie vaker gekozen wordt? 

Om deze vraag te beantwoorden maken we gebruik van een kansmodel. Door de context van het onderzoek en de eenvoud van de variabele (de keuze voor de helper of lastpost) kan dit kansmodel op intuïtieve wijze tot stand komen. De onderstaande lesactiviteit kan hier bij helpen. 
\end{multicols}

\clearpage

\begin{lesactiviteit}{Een kansmodel om `geen voorkeur' te simuleren}
	
\begin{vraagenantwoord}

    \vraag{Veronderstel even dat baby's (en dus ook de zestien baby's uit de studie) geen voorkeur hebben. Hoe vaak zal de helper gekozen worden? }
    \antwoord{Het exacte aantal kunnen we niet op voorhand voorspellen, maar we verwachten dat deze waarde in de buurt van acht zal liggen.}
		
    \vraag{Kun je een manier bedenken om dit proces (de keuze van de baby wanneer hij geen voorkeur heeft) na te bootsen?}
    \antwoord{We kunnen `geen voorkeur' simuleren door een geldstuk zestien keer op te werpen waarbij het gooien van kop staat voor het kiezen van de helper. Inderdaad, de kans op het kiezen van de helper (het werpen van kop) is gelijk aan de kans op het kiezen van de lastpost (het werpen van munt). Beide zijn vijftig procent.}
    
    \vraag{Vraag aan elke leerling om een geldstuk zestien keer op te werpen en het aantal keer dat kop gegooid werd neer te schrijven. Schrijf dit aantal op voor de verschillende leerlingen. Wat merk je? Indien je voldoende herhalingen hebt, kun je een staafdiagram aanmaken van deze reeks getallen.}
    \antwoord{Het aantal keer dat kop gegooid werd, varieert. Het ligt vaak in de buurt van acht, maar het kan hier ook van afwijken. }
    
    \vraag{Indien je het experiment heel vaak zou herhaald hebben, wat zou het staafdiagram dan voorstellen, gekoppeld aan de context van het onderzoek?}
    \antwoord{Het geeft de verwachte verdeling weer van het aantal baby's (op zestien) dat kiest voor de helper indien er geen voorkeur is. Of anders uitgedrukt: als de baby's geen voorkeur hebben en we herhalen de studie vele malen, dan geeft het staafdiagram de mogelijke waarden weer van het aantal keer dat de helper wordt gekozen. Hoe hoger de staaf, hoe meer waarschijnlijk de waarde is. }
    
    \vraag{Hoe waarschijnlijk is het dat veertien baby's of meer kiezen voor de helper indien ze geen voorkeur hebben?}
    \antwoord{Dit is onwaarschijnlijk. Als we zestien keer een geldstuk opwerpen, zullen we bijna altijd minder dan veertien keer kop gooien.}
\end{vraagenantwoord}
\end{lesactiviteit}

\clearpage

\begin{multicols}{2}

\autoref{fig:StaafdiagramKopBinomiaal} toont een staafdiagram van de experimentele kans bepaald door 720 keer 16 geldstukken op te werpen. Via een GeoGebra-app (Rollinson, 2023) kunnen we dit proces nabootsen (simuleren) zonder dat we daadwerkelijk geldstukken moeten opgooien. Net zoals in de lesactiviteit merken we dat veertien een resultaat is met een bijzonder kleine kans: bijna alle waarden zijn kleiner. Het is bijgevolg weinig waarschijnlijk om minstens veertien keer kop te werpen bij zestien worpen. 

\end{multicols}

\afbeelding{img/Lsimulatiepoppen.jpg}{Staafdiagram met de kansen voor het aantal keer kop bij 16 worpen met een geldstuk.}{fig:StaafdiagramKopBinomiaal}

\begin{multicols}{2}

Dit kunnen we vertalen naar de context van het onderzoek: veertien baby's die kiezen voor de helper is ongewoon veel indien ze geen voorkeur zouden hebben. We besluiten dat we bewijskracht hebben gevonden in de data dat het onwaarschijnlijk is dat de helper even vaak wordt gekozen als de lastpost (net zoals bij de oneerlijke dobbelstenen). Op dit moment is het nuttig om enkele begrippen te introduceren en die te koppelen aan het onderzoek. 

\begin{kader}{}{Begrippen}
\begin{itemize}
    \item \textbf{Nulhypothese $H_0$}. Een hypothese die we naar voor schuiven in de populatie en die we wensen te toetsen aan de hand van de data in de steekproef. Toegepast: $H_0$: de baby's hebben geen voorkeur. 
    \item \textbf{Toetsingsgrootheid}. Een maat die de data van de steekproef samenvat en toelaat de nulhypothese te evalueren. Toegepast: het aantal baby's dat kiest voor de helper. 
    \item \textbf{Nulverdeling}. De kansverdeling van de toetsingsgrootheid in de veronderstelling dat de nulhypothese waar is. Toegepast: de verdeling van het aantal keer dat de helper wordt gekozen als de baby's geen voorkeur hebben. 
\end{itemize}
\end{kader}

\columnbreak

Na deze eerste kennismaking met de binomiaaltoets gaan we enkele begrippen wat wiskundiger onderbouwen. Een sleutelrol is weggelegd voor de binomiale verdeling. 
	
    \section{De binomiale verdeling} \label{S_3}
     
    De \emph{binomiale verdeling} duikt op telkens wanneer een \emph{binair experiment} (ook wel \emph{Bernoulli experiment} genoemd) meermaals onafhankelijk en onder identieke omstandigheden uitgevoerd wordt. Met een binair experiment bedoelen we een experiment dat slechts twee uitkomsten kan hebben: ja of nee, kop of munt, besmet of niet besmet, wit of zwart, geslaagd of niet geslaagd ... In paragraaf \ref{S_2} ging het over het binaire experiment waarbij baby's kozen uit twee poppen: een helper of een lastpost.

    Wanneer we $n$ identieke binaire experimenten onafhankelijk van elkaar uitvoeren en we de kans op $i$ keer een bepaalde uitkomst berekenen (bijvoorbeeld de kans op $i$ keer ja, $i$ keer kop,  $i$ keer besmet, $i$ keer wit, $i$ keer geslaagd,  $i$ keer de helper ...) dan verkrijgen we per definitie een \emph{binomiale kansverdeling}. Doordat we hier spreken over identieke en onafhankelijke experimenten blijft voor elk binair experiment de kans op een specifieke uitkomst dezelfde. In het voorbeeld met de baby's wil dit zeggen dat de keuze van een bepaalde baby geen invloed kan hebben op die van een andere baby. De kans om de helper te kiezen wordt gelijk verondersteld bij elke baby.
    
    Soms moet de realiteit een beetje geforceerd worden om een binaire situatie en een perfecte binomiale verdeling te krijgen. Niet alles is zwart-wit. Niet elke persoon past bijvoorbeeld in de binaire man-vrouw-opsplitsing. Soms landt een muntstuk op zijn kant. Een binair experiment is dus vaak een vereenvoudigde afspiegeling van de werkelijkheid. Ook over de aanname van onafhankelijkheid valt het een en ander te zeggen.   
 

    \subsection*{Onafhankelijke experimenten}
	
    In de praktijk is het niet eenvoudig om de afhankelijkheid of de onafhankelijkheid van opeenvolgende experimenten aan te tonen. Als je via Tinder op zoek bent naar de liefde en iedere swipe als experiment beschouwt, is het duidelijk afhankelijk. Als je erg kieskeurig bent, is er slechts een kleine groep interessante kandidaten. Afhankelijk van je eerdere experimenten (swipes), zal de kans om een geschikte kandidaat te kunnen trekken dalen/stijgen. De (on)afhankelijkheid kan nog meer achter de schermen blijven. Als bijvoorbeeld leerlingen in een school tijdens een bevraging elkaars antwoorden kennen, kan er sprake zijn van afhankelijkheid door conformisme of door groepsdruk, zonder dat je daar als onderzoeker van op de hoogte bent.

    Laat ons terugkeren naar een wiskundiger voorbeeld met kauwgomballen.
    Als je kauwgomballen uit een automaat met rode en blauwe kauwgomballen laat komen en je vraagt je bij elke trekking af wat de kans is op een rode bal, dan merk je dat deze kans stilaan vermindert als je veel rode ballen na elkaar uit de automaat hebt laten rollen en dat ze vermeerdert wanneer je veel blauwe gekregen hebt. De kans op rood is bijgevolg niet constant. Dit effect is sneller merkbaar als het gaat over een kauwgomballenautomaat met een relatief kleine inhoud. Het aantal getrokken rode ballen leidt dan zeker niet tot een binomiale kansverdeling (maar wel tot een hypergeometrische kansverdeling). 
    
    Als je rode en blauwe kauwgomballen trekt uit een immense automaat, dan is de kans op rood nagenoeg constant en dan wordt de binomiale verdeling met dit herhaalde kansexperiment goed benaderd. Als je bovendien elke getrokken kauwgombal zou kunnen terugleggen na het bekijken van de kleur, dan is de kans op rood tijdens het trekken helemaal constant en krijg je dus een perfecte binomiale verdeling.

    \end{multicols}

    \afbeelding[10cm]{img/UWgrijpmachine.jpg}{}{}

    \begin{multicols}{2}

    \subsection*{Welke leerstof komt vooraf?}
    
    Het herkennen van een klokvormige grafiek als de dichtheidsfunctie van een normale verdeling is noodzakelijk voor deze paragraaf. De normale verdeling is sowieso erg belangrijk in de beschrijvende en in de verklarende statistiek. Daarom is het absoluut nodig dit leerstofonderdeel vooraf behandeld te hebben.

    In de vernieuwde tweede graad hebben de leerlingen geen kansbomen meer op het leerplan. De werkteksten hieronder maken gebruik van deze kansbomen, van de product- en de somregel en van de complementregel van kansen. We steunen bij de behandeling van de binomiale kansverdeling meermaals op deze basisprincipes. Deze leerstof moet dus ook op voorhand gezien zijn.

    
 \subsection*{Twee werkteksten}
 
    Voor de volgende werktekst hebben we het thema van de werkgelegenheidsgraad in België gekozen. Volgens StatBel (\url{statbel.fgov.be}) was $71,4\%$ van de bevolkingsklasse van 20 tot 64 jaar in België effectief aan het werk in 2021. Dit cijfer beweegt uiteraard in de tijd maar wij beschouwen het als constant en ronden het af naar 0,7. Hier hebben we de werkelijkheid vereenvoudigd en in een wiskundig keurslijf gekneed. 
	
\end{multicols}

\clearpage
\begin{lesactiviteit}{Werkgelegenheid en werkloosheid in België}

Op de website van StatBel, het Belgisch instituut voor statistiek, kun je opzoeken hoe groot de kans is dat een persoon uit de bevolkingsklasse van 20 tot 64 jaar effectief aan het werk is. Afgerond is deze kans 0,7. Aan de hand van dit getal kunnen we kanstheoretisch inschatten hoeveel personen in een steekproef van 100 personen er in deze klasse zitten en hoe groot de kans is dat deze schatting overschreden wordt. Zulke schattingen zijn van belang bij het toekennen van uitkeringen (werkloosheidssteun, ziekte-uitkering, studiebeurzen ...). In de volgende oefening bestuderen we nu eerst steekproeven van een kleinere omvang dan 100. We noemen de grootte van de steekproef $n$.

\begin{vraagenantwoord}
    \vraag{De kleinste steekproefgrootte die je kunt nemen is $n=1$. Stel $X$ de variabele die aangeeft hoeveel personen uit de steekproef een job hebben. Als $n=1$ zijn er slechts twee mogelijke uitkomsten: $X=0$ of $X=1$. Hoe groot is de kans op $X=0$ en de kans op $X=1$? Stel deze kansen voor met een staafdiagram.}
    
    \antwoord{Er zijn slechts twee staafjes: eentje voor $X=0$ en eentje voor $X=1$. De kans op $X=1$ is $0,7$. De kans op $X=0$ berekenen we met de complementregel: $1-0,7=0,3$. Hieronder zie je het gevraagde staafdiagram. We noemen dit staafdiagram een \emph{kanshistogram} omdat de oppervlakte van elk staafje een kans voorstelt. De variabele $X$ die deze kansen voortbrengt noemen we een \emph{kansvariabele}.}
    
    \afbeelding[6cm]{img/Lstaaf1.jpg}{}{} 
		
    \vraag{Bij een steekproefgrootte $n=2$ kies je ad random twee Belgen uit de bevolkingsklasse van 20 tot 64 jaar. Stel de kansvariabele $X$ opnieuw gelijk aan het aantal personen uit de steekproef met een job. Er zijn drie mogelijkheden: $X=0$, $X=1$ en $X=2$. De kansen op $X=0$, $X=1$ en $X=2$ noteren we wiskundig als $\textrm{P}(X=0)$, $\textrm{P}(X=1)$ en $\textrm{P}(X=2)$. De letter P is de eerste letter van \emph{probabilitas}, de Latijnse term voor kans. Teken een kansboom waarin de kansen op werkloosheid van twee willekeurige personen af te lezen zijn. Bereken hieruit $\textrm{P}(X=0)$, $\textrm{P}(X=1)$ en $\textrm{P}(X=2)$. Stel deze kansen daarna voor in een kanshistogram.}

    \antwoord{De kansboom heeft twee takken die beide in twee vertakkingen opsplitsen. Bij elke opsplitsing komen dezelfde kansen voor. Deze kansboom met vier takken helpt om de somregel en de productregel te visualiseren.
    \afbeelding[6cm]{img/Lkansboom.jpg}{}{}
    De kansen $\textrm{P}(X=0)$ en $\textrm{P}(X=2)$ berekenen we door de productregel toe te passen in de eerste en in de vierde tak van de kansboom. Maar bij $\textrm{P}(X=1)$ moeten we ook gebruik maken van de somregel. We moeten immers de tweede tak (waarbij persoon 1 werkloos is en persoon 2 een job heeft) combineren met de derde tak (waarbij persoon 1 een job heeft en persoon 2 werkloos is). Dit geeft de volgende kansen en het volgende kanshistogram.
    \begin{itemize}
	\item $\textrm{P}(X=0)=0,3 \cdot 0,3=0,09$
	\item $\textrm{P}(X=1)=0,3 \cdot 0,7+0,7 \cdot 0,3=2\cdot 0.21 =0,42$
	\item $\textrm{P}(X=2)=0,7 \cdot 0,7=0,49$
    \end{itemize}}
     
    \afbeelding[6cm]{img/Lstaaf2.jpg}{}{}
	
    \vraag{Los nu dezelfde vraag op voor $n=3$. Je kunt hiervoor een kansboom tekenen. Je mag de vraag ook oplossen zonder kansboom.}
    
    \antwoord{Als je een kansboom tekent, splits je de steekproeftrekking op in drie stappen met telkens twee vertakkingen. De kruin van de boom eindigt in acht takken. \\
    In plaats van op de takken van de kansboom te letten, kun je ook op de voegwoorden \emph{en} en \emph{of} letten. Voor $\textrm{P}(X=0)$ moet persoon 1 werkloos zijn, moet persoon 2 werkloos zijn \emph{en} moet persoon 3 werkloos zijn. Het voegwoord \emph{en} wijst op het gebruik van de productregel. Voor $\textrm{P}(X=3)$ moet persoon 1 een job hebben \emph{en} persoon 2 \emph{en} persoon 3. Dit leidt alweer tot een productregel. Bij $\textrm{P}(X=1)$ zit er precies één persoon met een job in de steekproef. De persoon met de job kan de eerste persoon zijn uit de steekproef \emph{of} de tweede \emph{of} de laatste. Het verbindingswoord \emph{of} wijst hier op een som (waarbij de termen geschreven worden als een product). \\ Door zowel kansbomen aan te bieden als te letten op voegwoorden kun je differentiëren tussen taalsterke en beeldsterke leerlingen. Beide redeneringen leiden tot dezelfde berekening en hetzelfde kanshistogram:
    \begin{itemize}
	\item $\textrm{P}(X=0)=0,3 \cdot 0,3 \cdot 0,3= 0,3^3=0,027$
        \item $\textrm{P}(X=1)=0,7 \cdot 0,3 \cdot 0,3+0,3 \cdot 0,7 \cdot 0,3+0,3 \cdot 0,3 \cdot 0,7=3\cdot 0,3^2\cdot 0,7 =0,189$
	\item $\textrm{P}(X=2)=0,7 \cdot 0,7 \cdot 0,3+0,7 \cdot 0,3 \cdot 0,7+0,3 \cdot 0,7 \cdot 0,7=3\cdot 0,3\cdot 0,7^2 =0,441$
	\item $\textrm{P}(X=3)=0,7 \cdot 0,7 \cdot 0,7= 0,7^3=0,343$
    \end{itemize}}
    
    \afbeelding[6cm]{img/Lstaaf3.jpg}{}{}

    \vraag{Er is een evolutie in de vorm van de kanshistogrammen: de eerste twee waren zuiver stijgende histogrammen maar het laatste histogram heeft ook een dalend verloop aan de rechterkant. Welke verdere evolutie merk je wanneer we verder gaan naar grotere steekproeven, bijvoorbeeld naar $n=10$ zoals op de onderstaande figuur?}
    \afbeelding[8cm]{img/Lstaaf6.jpg}{}{}
    \antwoord{Het aantal staafjes bij het stijgend verloop is meer in evenwicht met het aantal staafjes bij het dalend verloop. Het kanshistogram vertoont nog niet echt een symmetrie rond een verticale as.}

\end{vraagenantwoord}

Om de evolutie naar grote steekproeven nog beter te kunnen volgen, maken we nu een bruuske  overstap naar een steekproefgrootte $n=100$. 

\begin{vraagenantwoord}[resume]
    \vraag{Ga hiervoor naar GeoGebra en kies bij de drie bolletjes voor de optie \emph{kansrekening} of \emph{waarschijnlijkheidsrekening}. Het werkblad dat verschijnt, staat ingesteld op de normale verdeling. Wijzig dit in de binomiale verdeling. Kies voor $n=100$ en $p=0,7$. Bereken nu de kans dat het aantal personen met een job uit de steekproef van grootte 100 minstens 65 is en hoogstens 75, m.a.w. bereken $\textrm{P}(65 \le X \le 75)$. Bereken daarna de kans dat het aantal personen met een job minstens 73 is, m.a.w. bereken $\textrm{P}(73 \le X ).$}
    
    \antwoord{Na het ingeven van de eindpunten van de intervallen lezen we meteen af dat $\textrm{P}(65 \le X \le 75)$ gelijk is aan $0,7704$ en dat $\textrm{P}(73 \le X )$ gelijk is aan $0,2964$.}
	    
    \afbeelding[10cm]{img/Lstaaf4.jpg}{}{}
    \afbeelding[10cm]{img/Lstaaf5.jpg}{}{}
	
    \vraag{We merken dat de kanshistogrammen uit de vorige vraag klokvormig zijn. Op de trapjes na heeft de vorm veel weg van de normale verdeling. Verander nu de werkgelegenheidsgraad van 0,7 naar andere waarden. Blijft het binomiale kanshistogram dan klokvormig?}

    \antwoord{Ja, meestal wel. Maar in extreme gevallen is het diagram duidelijk links- of rechtsscheef en is het dus zeker niet meer klokvormig. Dat is bijvoorbeeld het geval als er zeer veel werklozen zijn (bijvoorbeeld bij $p=0,02$ en $n=100$) of wanneer bijna iedereen werkt (bijvoorbeeld bij $p=0,975$ en $n=100$). In GeoGebra is er een knop voorzien die de best passende klokvorm bovenop het kanshistogram tekent. Deze knop kan nuttig zijn bij het beantwoorden van deze vraag.}
	
    \vraag{Beantwoord nu de volgende vragen door binomiale kanshistogrammen in GeoGebra te tekenen.
    \begin{itemize}
        \item In België is 15\% van de bevolking rhesusnegatief. Hoe groot is de kans dat er minder dan 10 personen uit een aselecte steekproef van 50 personen rhesusnegatief zijn?
        \item Begin januari 2022 had 66\% van de inwoners van België een Belgische achtergrond. Hoe groot is de kans dat er bij een steekproef van 200 personen uit België meer dan 120 en minder dan 140 personen een Belgische achtergrond hebben? 
        \item Volgens StatBel heeft 73\% van de gezinnen in België een eigen wagen. Stel dat er een steekproef gedaan wordt van 73 Belgische gezinnen. Hoe groot is dan de kans dat er 60 of meer gezinnen een wagen bezit?
    \end{itemize}}
    
    \antwoord{Deze kansen zijn achtereenvolgens:
	\begin{itemize}
	    \item $\textrm{P}(X \le 9 )=0,7911$
	    \item $\textrm{P}(121 \le X \le 139 )=0,8247$
	    \item $\textrm{P}(60 \le X )=0,0465$
	\end{itemize}}
	
\end{vraagenantwoord}
\end{lesactiviteit}

\begin{multicols}{2}
    In deze loep kiezen we ervoor om het concept van de binomiale verdeling eerst goed te verankeren  vooraleer we overgaan naar de \emph{binomiaalcoëfficiënten}, die berekend worden met de formule voor het aantal \emph{combinaties} uit de \emph{combinatoriek}.  We leiden de formule voor de \emph{binomiaalcoëfficiënten} enkel en alleen af in functie van de binomiale verdeling. In de wiskundig georiënteerde studierichtingen past het begrip \emph{combinatie} in het bredere onderwerp van de telproblemen en wordt er naast combinaties ook aandacht besteed aan variaties en permutaties, al of niet met herhalingen. In de richting humane wetenschappen is dit bredere kader niet nodig. 
    
    De volgende werktekst sluit aan bij de voorgaande. We gaan nog even door met de werkgelegenheidskansen.
    
\end{multicols}

\clearpage
\begin{lesactiviteit}{Binomiaalcoëfficiënten en combinaties}   
    We hernemen één van de berekeningen uit de vorige werktekst. Als $X$ het aantal personen met een job is in een steekproef met $n=3$, dan vinden we door te tellen hoeveel takken van een kansboom we moeten samennemen onderstaande kansen:
    \begin{itemize}
        \item $\textrm{P}(X=0)=\textcolor{kleurloep}{1} \cdot 0,3 ^3$
	\item $\textrm{P}(X=1)=\textcolor{kleurloep}{3} \cdot 0,3^2 \cdot 0,7$
	\item $\textrm{P}(X=2)=\textcolor{kleurloep}{3} \cdot 0,3 \cdot 0,7^2$
	\item $\textrm{P}(X=3)=\textcolor{kleurloep}{1} \cdot 0,7^3$
    \end{itemize}

    De gekleurde getallen vooraan bij de berekening van de binomiale kansen noemen we de \emph{binomiaalcoëfficiënten}. Ze spelen een cruciale rol in kansberekening. In deze werktekst zoeken we eerst de binomiaalcoëfficiënten van een binomiale kansverdeling met $n=4$. Die vinden we makkelijk zonder specifieke formules in te voeren.

\begin{vraagenantwoord}
    \vraag{Hoe groot is de kans dat geen enkele persoon een job heeft in een steekproef van vier willekeurig gekozen Belgen uit de populatie van 20- tot 64-jarigen? En hoe groot is de kans dat ze alle vier een job hebben? Welke zijn de eerste en de laatste binomiaalcoëfficiënt bij $n=4$?}
    
    \antwoord{De eerste en de laatste binomiaalcoëfficiënt zijn gelijk aan \emph{1} want
    \begin{itemize}
        \item $\textrm{P}(X=0)=\textcolor{kleurloep}{1} \cdot 0,3 ^4$
	    \item $\textrm{P}(X=4)=\textcolor{kleurloep}{1} \cdot 0,7^4$
    \end{itemize}}
    
    \vraag{Hoe groot is de kans dat precies één persoon een job heeft? En hoe groot is de kans dat precies drie personen een job hebben. Welke zijn de tweede en de voorlaatste binomiaalcoëfficiënt in de rij van de binomiaalcoëfficiënten voor $n=4$?}

    \antwoord{De tweede en de voorlaatste binomiaalcoëfficiënt zijn gelijk aan \emph{4} want er zijn precies \emph{4} plaatsen waarop de enige \emph{0,7} (resp. de enige \emph{0,3}) kan staan in een product van \emph{4} factoren. Breder uitgeschreven: 
    \begin{itemize}
        \item $\textrm{P}(X=1)=0,7\cdot0,3\cdot0,3\cdot0.3+0,3\cdot0,7\cdot0,3\cdot0.3+0,3\cdot0,3\cdot0,7\cdot0.3+0,3\cdot0,3\cdot0,3\cdot0.7 = \textcolor{kleurloep}{4} \cdot 0,3 ^3 \cdot 0,7$
	    \item $\textrm{P}(X=3)=0,3\cdot0,7\cdot0,7\cdot0.7+0,7\cdot0,3\cdot0,7\cdot0.7+0,7\cdot0,7\cdot0,3\cdot0.7+0,7\cdot0,7\cdot0,7\cdot0.3 = \textcolor{kleurloep}{4} \cdot 0,3 \cdot 0,7^3$
    \end{itemize}}
    
    \vraag{Hoe groot is de kans dat precies twee van de vier aselect gekozen personen aan het werk zijn? Bepaal m.a.w. de middelste binomiaalcoëfficiënt in de rij van de binomiaalcoëfficiënten voor $n=4$.}
    
    \antwoord{De middelste binomiaalcoëfficiënt is \emph{6} want er zijn \emph{6} manieren om precies twee factoren \emph{0,7} te plaatsen in een product van \emph{4} factoren. De formule voor $\textrm{P}(X=2)$ breder uitgeschreven: \\ $\textrm{P}(X=2)=0,7\cdot0,7\cdot0,3\cdot0.3+0,7\cdot0,3\cdot0,7\cdot0.3+0,7\cdot0,3\cdot0,3\cdot0.7+0,3\cdot0,7\cdot0,7\cdot0.3+ +0,3\cdot0,7\cdot0,3\cdot0.7+0,3\cdot0,3\cdot0,7\cdot0.7= \textcolor{kleurloep}{6} \cdot 0,3 ^2 \cdot 0,7^2$}
\end{vraagenantwoord}

We vatten samen: de binomiaalcoëfficiënten voor $n=4$ zijn \textcolor{kleurloep}{1}, \textcolor{kleurloep}{4}, \textcolor{kleurloep}{6}, \textcolor{kleurloep}{4} en \textcolor{kleurloep}{1}.

We gaan nu over naar een steekproef met grootte $n=5$. Vanaf nu is het iets omslachtiger om de binomiaalcoëfficiënten op de klassieke manier te achterhalen. 

\lessubtitel{Klassieke berekening}

Stel dat je wilt berekenen hoe groot de kans is dat precies twee van de vijf willekeurig gekozen Belgen een job hebben. Als je de steekproefpersonen met een job met de letter J voorstelt en de werklozen met een W dan kun je snel een overzicht maken van alle rijtjes van vijf personen: twee met een job en drie zonder.

\begin{vraagenantwoord}[resume]

\vraag{Doe dit op een systematische manier.}
\antwoord{Hieronder zie je een van de mogelijkheden.
    \[
    \begin{array}{c|c|c|c|c}
        \cellcolor{kleurloep}{J} & \cellcolor{kleurloep}{J} & W & W & W\\
        \hline
        \cellcolor{kleurloep}{J} & W & \cellcolor{kleurloep}{J} & W & W\\
        \hline
        \cellcolor{kleurloep}{J} & W & W & \cellcolor{kleurloep}{J} & W\\
        \hline
        \cellcolor{kleurloep}{J} & W & W & W & \cellcolor{kleurloep}{J}\\
        \hline
        W & \cellcolor{kleurloep}{J} & \cellcolor{kleurloep}{J} & W & W
    \end{array}
    \qquad
    \begin{array}{c|c|c|c|c}
        W & \cellcolor{kleurloep}{J} & W & \cellcolor{kleurloep}{J} & W\\
        \hline
        W & \cellcolor{kleurloep}{J} & W & W & \cellcolor{kleurloep}{J}\\
        \hline
        W & W & \cellcolor{kleurloep}{J} & \cellcolor{kleurloep}{J} & W\\
        \hline
        W & W & \cellcolor{kleurloep}{J} & W & \cellcolor{kleurloep}{J}\\
        \hline
        W & W & W & \cellcolor{kleurloep}{J} & \cellcolor{kleurloep}{J}
    \end{array}
    \]
}

\end{vraagenantwoord}

In totaal zijn er 10 van deze rijtjes. We noteren dit aantal als $C^2_5=10$ waarbij $C$ de eerste letter is van het woord \emph{combinatie}. De notatie  $C^2_5$ gebruiken we in het algemeen om het \emph{aantal combinaties} van 2 uit 5 aan te geven. In deze concrete toepassing staat het getal 2 bovenaan voor het aantal personen met een job (of het aantal gekleurde hokjes). Het getal 5 onderaan staat hier voor de steekproefgrootte $n$ (of het totale aantal hokjes). 

Op deze manier hebben we nu op een omslachtige manier een nieuwe binomiaalcoëfficiënt gevonden. Hiermee kunnen we de kans op twee personen met een job uit een steekproef van vijf personen berekenen: $\textrm{P}(X=2)=\textcolor{kleurloep}{10} \cdot 0,3^3 \cdot 0,7^2.$

\lessubtitel{Snellere berekening}

Sneller gaat het als je beredeneert op hoeveel manieren je twee van de vijf vierkantjes op een rij kunt inkleuren. Er zijn 5 mogelijke vierkantjes die je als eerste kunt inkleuren. Er blijven nog 4 vierkantjes die je hierna kunt inkleuren. Maar het totale aantal inkleuringen is toch niet gelijk aan $5 \cdot 4$. 

\begin{vraagenantwoord}[resume]
    \vraag{Verklaar dit. Hoe bereken je $C^2_5$ dan wel? Gebruik daarna de waarde van $C^2_5$ om de kans te berekenen op 2 personen met een job in een steekproef van 5 stuks?}
    \antwoord{Als eerst het derde vierkantje wordt ingekleurd en daarna het vijfde dan heeft dit hetzelfde effect als wanneer eerst het vijfde vierkantje wordt ingekleurd en daarna het derde. Elke kleuring wordt hierboven dus twee keer geteld. Omdat de volgorde van geen belang is, moeten we dus nog delen door twee: $C^2_5=\dfrac{5 \cdot 4}{2}=10$. Met deze binomiaalcoëfficiënt vinden we de kans op 2 personen met een job in een steekproef van 5 personen: $\textrm{P}(X=2)=\textcolor{kleurloep}{10} \cdot 0,3^3 \cdot 0,7^2.$}
\end{vraagenantwoord}

\clearpage

Iets moeilijker nu. Je neemt een steekproef van zeven personen en je wilt de kans kennen op drie personen met een job. Om de binomiaalcoëfficiënt te kennen, bereken je op hoeveel manieren je drie vierkantjes kunt inkleuren in een rijtje van zeven. Voor het eerste vierkantje zijn er 7 mogelijkheden, voor het tweede zijn er 6 en voor het derde resten er 5 mogelijkheden. Toch is het aantal inkleuringen niet gelijk aan $7 \cdot 6 \cdot 5$.

\begin{vraagenantwoord}[resume]
    \vraag{Verklaar dit. Hoe bereken je $C^3_7$ dan wel? En hoe groot is de kans dan op 3 personen met een job in een steekproef van 7 stuks?}
    \antwoord{Als de drie vierkantjes in een andere volgorde worden ingekleurd dan ziet het resultaat er identiek uit. We gaan dus na op hoeveel manieren drie vaste vierkantjes kunnen ingekleurd worden. Voor het eerste van die vierkantjes zijn er 3 mogelijkheden, voor het tweede 2 en voor het laatste 1. De drie vaste vierkantjes kunnen dus op $3 \cdot 2 \cdot 1$ manieren in een verschillende volgorde ingekleurd worden, 6 manieren dus. We vinden dus dat $C^3_7=\dfrac{7 \cdot 6 \cdot 5}{1 \cdot 2 \cdot 3}=35$. Met deze binomiaalcoëfficiënt vinden we de kans op 3 personen met een job in een steekproef van 7 personen: $\textrm{P}(X=3)=\textcolor{kleurloep}{35} \cdot 0,3^4 \cdot 0,7^3.$ }
    
     \vraag{Hoe groot is de kans op 3 gezinnen met een wagen in een steekproef van 7 stuks? Je mag nog altijd veronderstellen dat 73\% van de Belgische gezinnen over een wagen beschikt.}
    \antwoord{ $\textrm{P}(X=3)=\textcolor{kleurloep}{35} \cdot 0,27^4 \cdot 0,73^3.$ }
    
    \vraag{Bereken $C^4_9$. Stel hiermee een formule op voor de kans dat je 4 keer een 3 gooit als je 9 keer met een gewone dobbelsteen gooit.}
    \antwoord{Door te tellen op hoeveel manieren we vier vierkantjes kunnen inkleuren in een rijtje van 9 vinden we $C^4_9=\dfrac{9 \cdot 8 \cdot 7 \cdot 6}{1 \cdot 2 \cdot 3 \cdot 4}=126.$ Verder hebben we nog de kans nodig op een 3-worp. Die kans is gelijk aan $\dfrac{1}{6}$. De kans op een ander aantal ogen dan 3 is $\dfrac{5}{6}$. Hieruit berekenen we dan dat $\textrm{P}(X=4)=\textcolor{kleurloep}{126} \cdot \left( \dfrac{1}{6} \right) ^4 \cdot \left( \dfrac{5}{6} \right) ^5$. De kans op 4 keer een 3 te gooien in 9 beurten is iets minder 4\%.}
    
\end{vraagenantwoord}
\end{lesactiviteit}

\begin{multicols}{2}

Na deze tweede werktekst is het duidelijk hoe de kansen in een binomiaal experiment kunnen berekend worden door de kennis van de binomiaalcoëfficiënten en dat de binomiaalcoëfficiënten kunnen berekend worden met een formule voor het aantal combinaties van $k$ vierkantjes uit een reeks van $n$ vierkantjes. 

Een eerste formule voor aantallen combinaties maakt gebruik van producten met puntjes in het midden (ook wel een \emph{gedurig product} genoemd), zie (5.1). In de teller staan vanaf $n$ afdalende factoren en in de noemer vanaf $1$ opklimmende factoren. Het aantal factoren in de teller (en in de noemer) is $k$.

Een tweede formule maakt gebruik van de faculteitsfunctie, zie (5.2). Hoewel de faculteitsfunctie ook een product met puntjes voorstelt, geniet deze de voorkeur als er met software gewerkt wordt en de waarde van $n$ groot is. Door het gebruik van de faculteitsfunctie $n$! $=1 \cdot 2 \cdot 3 \cdot \ldots \cdot n$ kan het typwerk heel wat ingekort worden. Maar als het werken met de faculteitsnotatie een stap te ver is voor de leerlingen, kan die ook weggelaten worden. Bij het gedurige product zien de leerlingen duidelijker wat er gebeurt.

Het verband tussen beide formules kan in de klas aangetoond worden via een voorbeeld. 

\begin{kader}{Formule}{Aantal combinaties berekenen met een product met puntjes}
\begin{equation}
    C^k_n=\frac{n \cdot (n-1) \cdot \ldots \cdot (n-k+1)}{1 \cdot 2 \cdot \ldots \cdot k}
    \begin{array}{l}
        \rightarrow k \text{ factoren}\\
        \rightarrow k \text{ factoren}
    \end{array}
\end{equation}
\end{kader}

\clearpage

\begin{kader}{Formule}{Aantal combinaties berekenen met faculteiten}
\begin{equation}
    C^k_n=\frac{n!}{k! \cdot (n-k)!}
\end{equation}
\end{kader}

 We merken hier nog even op dat het werken met faculteiten op een zakrekenmachine beperkingen heeft. Op mijn TI84 Plus worden de resultaten bij een faculteitsberekening te groot vanaf $70!$. Er verschijnt een overflow-error. Echter, bij de berekening van het getal $C^{32}_{70}$ via het menu Kansrekening ontstaat deze overflow-error niet. Het algoritme voor de berekening van aantallen combinaties vereenvoudigt tellers en noemers wellicht tijdig.
 
 De GRexit (zoals besproken in Uitwiskeling 38/3) zorgt er weliswaar voor dat we meer rekenkracht ter beschikking hebben. Maar de opmerking blijft gelden. Ook computerapps zoals Calculator-$\infty$ hebben geen enkel probleem met $C^{70}_{180}$ ook al zijn ze niet in staat om 180! te berekenen.

 Tot slot wordt de algemene formule voor de binomiale kansverdeling opgesteld. Bij een binomiaal experiment met een succeskans $p$ en een mislukkingskans $1-p$ wordt de kans $\textrm{P}(X=k)$ op $k$ successen en $n-k$ mislukkingen gegeven door de volgende formule. 

 \begin{kader}{Formule}{Kansen van een binomiale stochast}
\begin{equation}
\begin{split}
&\textrm{P}(X=k) = C^k_n p^{k} (1-p)^{n-k} \\
&\textrm{met }k=0,1,\ldots,n
\label{eq:Eqbinomiale}
\end{split}
\end{equation}
\end{kader}
 
\subsection*{Het gemiddelde van een binomiaal verdeelde veranderlijke}

Deze paragraaf houden we bewust heel kort. We brengen de formule voor het gemiddelde aan met een voorbeeld. 

Stel dat in een bepaalde school de kans op een leerling met een thuistaal die niet het Nederlands is, gelijk is aan $p=\frac{1}{4}$. Het aantal leerlingen in het vijfde jaar van deze school is $n=48$. Van hoeveel van deze leerlingen mag je dan verwachten dat ze thuis niet het Nederlands als voertaal gebruiken? Anders geformuleerd: hoeveel leerlingen zullen er gemiddeld gezien niet het Nederlands als thuistaal gebruiken in vergelijkbare jaren (met $n=48$) van vergelijkbare scholen (met $p=\frac{1}{4}$)? 

Hier zal het voor de leerlingen meteen duidelijk zijn dat er ongeveer $48 \cdot \frac{1}{4}$ of 12 anderstalige leerlingen zullen zijn. De algemene formule voor de verwachtingswaarde van een binomiaal verdeelde veranderlijke ligt nu voor de hand:

\[E(X)=n\cdot p.\]

Op een grafiek zien we dat dit gemiddelde op de horizontale as precies onder de top van de grafiek ligt, zie \autoref{fig:gemiddelde}.

\afbeelding[8cm]{img/Lgemiddelde.jpg}{De verwachtingswaarde ligt onder de top}{fig:gemiddelde}

Eventueel kunnen we nu ook iets vertellen over de spreiding van de gegevens. Als de kans op een leerling met een andere thuistaal heel klein zou zijn ($p$ dicht bij 0) of net zeer groot ($p$ in de buurt van 1) dan is er geen grote spreiding meer van het aantal anderstaligen. In het ene geval zal het aantal leerlingen met een andere thuistaal voornamelijk in een kleine zone dicht bij 0 liggen en in het andere geval in een kleine zone dicht bij 48. Een goede maatstaf voor de standaardafwijking $\sigma(X)$ van de binomiale kansverdeling wordt gegeven door de formule

\[\sigma(X)=\sqrt{n\cdot p\cdot(1-p)}\]

want voor $p=0$ en voor $p=1$ bereikt deze formule een minimum. 

In de context van de leerlingen die niet het Nederlands als thuistaal hebben, kunnen we narekenen dat $\sigma(X)=\sqrt{9}=3$. Als we vanuit het gemiddelde een standaardafwijking naar links en een standaardafwijking naar rechts gaan, verkrijgen we de grenzen van het interval [9,15] zoals in \autoref{fig:standaard}.

\afbeelding[8cm]{img/Lstandaardafwijking.jpg}{Vanuit het gemiddelde een standaardafwijking naar links en een standaardafwijking naar rechts}{fig:standaard}

Met deze waarden kun je een link leggen naar de normale verdeling, die in dit voorbeeld goed aansluit bij de klokvormige binomiaalverdeling. De waarden 9 en 15 op de horizontale as, zijn de waarden waarin de grafiek overgaat van hol naar bol of van bol naar hol. We noemen deze punten de buigpunten.

\afbeelding[8cm]{img/Lstandaardafwijking2.jpg}{De normaalverdeling die best aansluit bij de binomiale verdeling}{fig:standaard2}

Nu we voldoende kennis hebben over de binomiale verdeling kunnen we terugkeren naar de binomiaaltoets waarin we deze verdeling voor het eerst nodig hadden.

\section{Toetsen van hypotheses (vervolg)} \label{S_4}
In deze paragraaf beschrijven we formeel de vijf stappen in het toetsen van hypotheses. We keren daarvoor terug naar het onderzoek rond moraliteitsbesef bij baby's en zullen daarbij de binomiale verdeling uit de vorige paragraaf gebruiken. Hoewel de berekeningen bij andere hypothesetoetsen gebruik kunnen maken van andere verdelingen, zoals de normale verdeling of student-t-verdeling, zijn de vijf besproken stappen algemeen geldig.
%de binomiaaltoets formeel in vijf stappen en keren we terug naar het onderzoek rond moraliteitsbesef bij baby's. 

\subsection*{Stap 1: hypotheses formuleren}

Bij een hypothesetoets starten we met het formuleren van de hypotheses (de veronderstellingen). Meer specifiek zullen er twee hypotheses zijn: de \emph{nulhypothese} \(H_0\) en de \emph{alternatieve hypothese} \(H_A\). De formulering van de hypotheses gebeurt op basis van de onderzoeksvraag waarbij de nulhypothese typisch stelt dat er geen effect is op populatieniveau (bv. er is geen effect van een testmedicijn op depressie of er is geen associatie tussen roken en longkanker) en de alternatieve hypothese stelt dat er wel een effect is op populatieniveau (bv. er is een effect van een testmedicijn op depressie of er is een associatie tussen roken en longkanker). 

Het onderzoek beoogt dan op basis van de verzamelde data bewijskracht te vinden tegen de nulhypothese en voor de alternatieve hypothese. Toegepast op het onderzoek rond moraliteitsbesef bij baby's wordt dit \(H_0:\) de helper wordt even vaak gekozen als de lastpost (er is \emph{geen} voorkeur en dus geen moraliteitsbesef bij baby's) en \(H_A:\) de helper wordt vaker gekozen dan de lastpost (er is moraliteitsbesef bij baby's).  

We stellen de voorkeur van een baby \(i\) voor door de variabele \(Y_i\), waarbij \(Y_i=1\) staat voor het kiezen van de helper en \(Y_i=0\) voor het kiezen van de lastpost. We hebben hier de twee waarden (helper of lastpost) numeriek gecodeerd. Dit zal handig zijn als we straks de absolute frequentie compact willen neerschrijven.

Door \(\textrm{P}(Y_i = 1)\) stellen we de kans voor dat baby \(i\) de helper kiest. Deze kans zullen we beknopt schrijven via \(p\), dus \(p = \textrm{P}(Y_i = 1)\), en gelijk veronderstellen voor elke baby. Of in woorden: \(p\) is de kans dat een baby de helper kiest.

We kunnen de hypotheses compact schrijven als
\[
H_0: p = 0,5 \quad \text{versus} \quad H_A:  p > 0,5. 
\]
Met andere woorden stelt de nulhypothese dat de kans dat de helper wordt gekozen gelijk is aan 50\% terwijl de alternatieve hypothese stelt dat deze kans groter is dan 50\% (de helper wordt vaker gekozen dan de lastpost). Merk op dat we bij de alternatieve hypothese niet aangeven hoeveel groter de kans dan exact is.

\subsection*{Stap 2: toetsingsgrootheid vastleggen}
In een volgende stap defini"eren we een \emph{toetsingsgrootheid} \(T\). De toetsingsgrootheid vat informatie uit de data samen op een manier die aansluit bij de hypothese die je wenst te toetsen. Omdat de hypotheses geformuleerd zijn in termen van de kans dat een baby kiest voor helper, vormt het aantal baby's dat de helper heeft gekozen in de steekproef een natuurlijke keuze. Dankzij de numerieke codering van \(Y_i\) kunnen we voor een steekproef van grootte $n$ de toetsingsgrootheid compact neerschrijven als:
\[
T = \sum_{i=1}^n Y_i. 
\]
We schrijven de toetsingsgrootheid met een hoofdletter \(T\) omdat ze een variabele is. De absolute frequentie kan inderdaad wijzigen als we de studie herhalen, m.a.w. ze kan \emph{vari\"eren} van steekproef tot steekproef. Als we de toetsingsgrootheid berekenen voor de data geobserveerd in de specifieke steekproef, dan noemen we ze de \emph{geobserveerde toetsingsgrootheid} en schrijven we ze als \(t_{obs}\). Toegepast op onze studie wordt dit \(t_{obs}=14\).

\subsection*{Stap 3: het kansmodel van de toetsingsgrootheid bepalen}


In deze stap  willen we de \emph{steekproevenverdeling} van de toetsingsgrootheid \(T\) bepalen. We bedoelen hiermee: welke waarden kan \(T\) aannemen indien we de studie vele malen zouden herhalen en wat is de kans op deze waarden? De steekproevenverdeling geeft hierop het antwoord. Bij het opstellen van de steekproevenverdeling zien we data als de uitkomst van een toevalsproces en het is dit proces dat we via een \emph{kansmodel} zullen modelleren. Meer specifiek gebruiken we de binomiale verdeling als kansmodel. 

De binomiale verdeling zal een goede benadering vormen van het toevalsproces indien we enkele veronderstellingen maken:

\begin{itemize}
\item
  De kans dat een baby de helper kiest, wordt constant verondersteld bij elke trekking uit de steekproef.
\item
  De keuzes van de baby's zijn onafhankelijk van elkaar in de steekproef.
 \item 
 Het aantal baby's dat deelneemt aan de studie (i.e. de grootte van de steekproef) ligt op voorhand vast. 
\end{itemize}
De eerste veronderstelling kunnen we niet makkelijk nagaan. Misschien hangt de kans af van bepaalde karakteristieken van de baby. Gelukkig is deze assumptie niet cruciaal. Zelfs wanneer de kansen niet constant zijn, zal de binomiale verdeling een goede benadering geven van het toevalsproces. 

Aan deze tweede veronderstelling kan voldaan zijn omdat de baby's het poppenspel afzonderlijk te zien te krijgen. We kunnen echter geen volledige zekerheid hebben of aan deze veronderstelling voldaan is. Als bijvoorbeeld de helft van de baby's het poppenspel te zien krijgen door een eerste proefleider en de andere helft door een tweede proefleider en als de proefleider een onbewuste invloed heeft op de keuze, dan is er niet voldaan aan deze veronderstelling. 

De derde veronderstelling is de makkelijkste om aan te voldoen door op voorhand de steekproefgrootte vast te leggen. 

Veel statistische technieken doen beroep op assumpties (of veronderstellingen) en vaak zal het niet eenvoudig zijn om te `bewijzen' dat de assumpties voldaan zijn. Het is daarom belangrijk om steeds de assumpties te rapporteren. Als de assumpties niet voldaan zouden zijn, stelt ons kansmodel een vereenvoudiging van de werkelijkheid voor. We hopen dan dat deze vereenvoudiging toch nog waardevolle inzichten oplevert.


Als de veronderstellingen voldaan zijn, kun je de mogelijke waarden van \(T\) en de kansen waarmee deze voorkomen, bij een herhaling van de studie onder dezelfde condities modelleren via de \emph{binomiale verdeling} die we hierboven reeds uitvoerig hebben besproken.

%\begin{kader}{Formule}{Kansen van een binomiale stochast}
%\begin{equation}
%\begin{split}
%&P(T = k) = \frac{n!}{k!(n-k)!} p^{k} (1-p)^{n-k} \\
%&\textrm{met }k=0,1,\ldots,n
%\label{eq:Eqbinomiale}
%\end{split}
%\end{equation}
%\end{kader}

Als we de uitdrukking van de binomiale verdeling \eqref{eq:Eqbinomiale} willen gebruiken om de kans te berekenen dat $k$ van de $n$ baby's de helper kiezen, moeten we weten welke waarde \(p \) aanneemt. We moeten dus de kans weten dat een baby kiest voor de helper. Deze kans kennen we niet en we voeren net de hypothesetoets uit om een uitspraak te kunnen doen over deze kans. De binomiale verdeling geeft ons dus een wiskundige formule voor de steekproevenverdeling, maar we kunnen ze helaas niet direct gebruiken omdat \(p\) ongekend is. De nulverdeling zal ons uit deze impasse helpen.

\subsection*{Stap 4: de nulverdeling opstellen}

De \emph{nulverdeling} is  de steekproevenverdeling van \(T\) wanneer we veronderstellen dat de nulhypothese opgaat. Als de nulhypothese waar is dan is \(p =0.5\) en kunnen we dit invullen in de uitdrukking van de binomiaalverdeling \eqref{eq:Eqbinomiale}. \autoref{fig:nulverdeling} visualiseert deze verdeling onder de nulhypothese. 

\afbeelding[8cm]{img/LNulverdeling.pdf}{De nulverdeling}{fig:nulverdeling}



\subsection*{Stap 5: de P-waarde berekenen en interpreteren}

We weten al dat de nulverdeling de verdeling van de toetsingsgrootheid weergeeft als de nulhypothese waar is. Of anders uitgedrukt: de nulverdeling geeft aan met welke kans het toevalsproces bepaalde waarden kan aannemen als de nulhypothese waar is. Of in dit specifieke voorbeeld: de nulverdeling geeft de kansen weer wanneer we de studie zouden herhalen, van hoeveel baby's er zouden kiezen voor de helper indien ze geen voorkeur hebben.

Als de alternatieve hypothese opgaat - dus als de kans om de helper te kiezen groter is dan 50\% - dan verwachten we dat er meer baby's zullen kiezen voor de helper dan wanneer de nulhypothese opgaat. Wanneer de alternatieve hypothese waar zou zijn, verwachten we dus toetsingsgrootheden die groter zullen zijn dan wanneer de nulhypothese waar zou zijn. Dit stelt ons nu in staat om bewijskracht in de data tegen de nulhypothese en in het voordeel van de alternatieve hypothese uit te drukken: we berekenen de kans om een waarde te bekomen groter dan of gelijk aan de geobserveerde toetsingsgrootheid in de veronderstelling dat de nulhypothese waar is. Deze kans wordt de \emph{P-waarde} genoemd. 

Formeel kunnen we de P-waarde schrijven als \(\textrm{P}(T \geq t_{obs} \mid H_0)\) waarbij je de notatie \(\textrm{P}(A \mid H_0)\) leest als de kans op \(A\) als \(H_0\) waar is. Toegepast op het voorbeeld waarbij \(t_{obs} = 14\) wordt dit \(\textrm{P}(T \geq 14 \mid H_0)\), de kans dat de toetsingsgrootheid minstens 14 bedraagt indien de nulhypothese waar is.

We kunnen de P-waarde ook visualiseren op de nulverdeling in \autoref{fig:nulverdeling}. Ze is gelijk aan de som van de oppervlaktes van de staven die horen bij de toetsingsgrootheden 14, 15 en 16. Als we deze kans berekenen in GeoGebra krijgen we $\textrm{P}(T \geq 14 \mid H_0) = 0,0021$. Indien de nulhypothese waar is, dus indien de baby's geen voorkeur vertonen, dan verwachten we bijna nooit dat 14 of meer baby's zullen kiezen voor de helper. Omdat deze kans zeer klein is, kunnen we stellen dat de geobserveerde toetsingsgrootheid ongewoon groot is om zich voor te doen als de nulhypothese waar is. 
%(modus tollens en inductie).

De P-waarde kwantificeert de bewijskracht in de steekproef tegen de nulhypothese: hoe kleiner de P-waarde, hoe minder waarschijnlijk het is dat de geobserveerde steekproefdata zich voordoen als de nulhypothese waar is. Hoe kleiner de P-waarde, hoe meer bewijskracht we dus hebben tegen de nulhypothese. Of anders gezegd: hoe kleiner de P-waarde, hoe minder plausibel de nulhypothese is in het licht van de geobserveerde data en de alternatieve hypothese. 

Wanneer in ons voorbeeld veertien van de zestien baby's kiezen voor de helper, concluderen we dus dat het weinig plausibel is dat de baby's geen voorkeur hebben. We verwerpen m.a.w. de nulhypothese in het voordeel van de alternatieve hypothese. Merk op dat we niet zeggen dat de nulhypothese vals is, noch dat de alternatieve hypothese waar is. We zitten immers nog steeds met onzekerheid. Dit maakt dat statistische besluitvorming verschilt van de klassieke logische besluitvorming met bijvoorbeeld modus tollens. 

De hypothesetoets laat ons  dus toe om de bewijskracht in de data tegen de nulhypothese en in het voordeel van de alternatieve hypothese uit te drukken. Deze bewijskracht kunnen we nu gebruiken om een besluit te nemen: verwerpen we de nulhypothese op basis van de data of verwerpen we ze niet? De beslissingsregels die we hanteren doen beroep op een \emph{significantieniveau} dat we noteren met $\alpha$. Dit significantieniveau leggen we vast vooraleer we het onderzoek uitvoeren. Dit getal is vaak klein. In de praktijk gebruikt men vaak $\alpha=0,05$, maar andere keuzes zijn ook mogelijk. De beslissingsregels zijn dan als volgt:

\begin{itemize}
\item Indien P-waarde  $\leq \alpha$: verwerp de nulhypothese. We beslissen dat we in de data voldoende bewijs gevonden hebben tegen de nulhypothese.
\item Indien P-waarde  $> \alpha$: verwerp de nulhypothese niet. We beslissen dat we in de data onvoldoende bewijs hebben gevonden tegen de nulhypothese.
\end{itemize}

Wanneer we de nulhypothese verwerpen, spreken we over een \emph{statistisch significant} resultaat. 

\subsection*{Eenzijdige en tweezijdige toetsen}
In het bovenstaande voorbeeld formuleerden we de alternatieve hypothese als $H_A:  p > 0,5$. Dit is een voorbeeld van een \emph{eenzijdige hypothese}: we zoeken enkel bewijs tegen $H_0$ dat erop wijst dat de helper vaker gekozen wordt. We kijken dus maar naar afwijking van $H_0$ in één richting: de richting waar de helper vaker wordt gekozen. Dit is de richting die we zouden verwachten volgens het moraliteitsprincipe.

In de praktijk zal men eenzijdige alternatieve hypotheses maar gebruiken indien men \emph{a priori} vrij zeker is wat de juiste richting is. Deze keuze mag je dus niet laten afhangen van de data. In de empirische cyclus moeten we immers eerst de hypotheses formuleren, vervolgens data verzamelen en die dan gebruiken om de voorgestelde hypothese te bevestigen of te ontkrachten. 

Wanneer men dus op voorhand geen idee heeft over de richting van het alternatief, dan zal men opteren voor \emph{de tweezijdige alternatieve hypothese}, bv. $H_A: p \neq 0,5$, dat bewijs zoekt in de data tegen de nulhypothese en dit in beide richtingen. De tweezijdige alternatieve hypothese kun je dus zien als een veilige keuze. Ze wordt in de praktijk vaak gebruikt. 

\subsection*{Type-I- en type-II-fouten}

Afhankelijk van de vragen van je leerlingen of de vereiste diepgang in de (nog te verschijnen) eindtermen kan het nuttig zijn om ook in te gaan op de mogelijke fouten die bestaan bij hypothesetoetsen. We houden hierbij de visie aan dat het beter is om te focussen op de praktische interpretatie hiervan dan op de grafische betekenis als oppervlakte onder een dichtheidskromme. Zo lijkt het ons wel relevant dat leerlingen beseffen dat deze fouten inherent zijn aan het toetsingsproces maar niet relevant om de berekening ervan expliciet uit te voeren. De volgende lesactiviteit geeft een mogelijke aanpak. We maken hierbij gebruik van een metafoor met de rechtspraak.

\end{multicols}

\begin{lesactiviteit}{Type-I- en type-II-fouten bij rechtspraak}
\begin{vraagenantwoord}
        \vraag{Bij hypothesetoetsen werken we steeds met twee uitspraken waarbij we bewijskracht verzamelen tegen de ene hypothese in het voordeel van de andere hypothese. Binnen de rechtspraak kennen we een heel gelijkaardige situatie. Weet je welke?}
        \antwoord{Binnen een rechtzaak gaat men ook steeds na of er voldoende aanwijzingen zijn om de schuld van een verdachte vast te stellen. Indien dit niet het geval is, dan wordt de aanname van onschuld behouden.}
        
        \vraag{De jury of rechter moet op het einde dus steeds een beslissing maken over de schuld of onschuld van de beklaagde. Welke twee
        foute uitspraken zijn er mogelijk?}
        \antwoord{Een onschuldige veroordelen of een schuldige vrijspreken.}
         
        \vraag{Is het mogelijk om beide fouten te vermijden?}
        \antwoord{Neen: om de eerste fout (een onschuldige veroordelen) uit te sluiten moet je iedereen vrijspreken en zo maak je net veel fouten van de tweede soort (een schuldige vrijspreken). Anderzijds: om de tweede fout uit te sluiten moet je iedereen veroordelen en zo maak je net veel fouten van de eerste soort. De fouten zijn dus inherent aan het beslissingsproces.}
        
        \vraag{Welke fout is volgens jullie het ergste?}
        \antwoord{Meestal vindt het merendeel van de klas dat het veroordelen van een onschuldige erger is.}
        
        \vraag{Het Belgisch grondbeginsel van het strafrecht neemt een duidelijk standpunt in betreffende de vorige vraag. Weet je wat dit beginsel zegt?}
        \antwoord{Je bent onschuldig tot het tegendeel bewezen is. Dit principe volgt zelfs uit de universele verklaring van de rechten van de mens.}
        \end{vraagenantwoord}
        \end{lesactiviteit}

\begin{multicols}{2}
 De leerlingen kunnen het strafrecht dus zien als een metafoor voor een hypothesetoets.
 
 De nulhypothese is in de rechtspraak het vermoeden van onschuld. 
 Ook in de statistiek staat de nulhypothese vaak voor een bewering waarvan we realistisch kunnen vertrekken of die algemeen geaccepteerd wordt bij gebrek aan bewijs dat de bewering zou tegenspreken. Bijvoorbeeld: baby's kiezen even vaak voor het ene poppetje als voor het andere. Bij een medische hypothesetoets staat de nulhypothese meestal voor de bewering die stelt dat er gemiddeld geen merkbaar verschil is tussen pati\"enten die een nieuwe, onbekende behandeling krijgen en diegenen die ze niet krijgen.
%Ook in de statistiek staat de nulhypothese vaak voor een bewering die stelt dat een zekere behandeling geen effect heeft, er geen associatie is, \ldots

 De alternatieve hypothese is in de rechtspraak de reden waarom een verdachte moet verschijnen voor de rechtbank. De aanklager wenst deze hypothese aan te tonen. Ook in de statistiek is de alternatieve hypothese degene die men wenst aan te tonen. In het bovenstaande medisch voorbeeld zou de alternatieve hypothese overeenkomen met een merkbaar, positief effect van de nieuwe behandeling.

 Zowel in de rechtspraak als in de statistiek is geen absoluut bewijs mogelijk en moet er dus een standaardprincipe vooropgesteld worden om de nulhypothese te kunnen verwerpen. In de rechtspraak moet er ook `voldoende bewijs' zijn om iemand te kunnen veroordelen. We gebruiken  steeds maar parti\"ele informatie/data omdat het onmogelijk is om de volledige waarheid te kennen in de echte wereld.
 
 Hierdoor kunnen er twee soorten van fouten ontstaan:
 \begin{itemize}
\item Type I: onterecht een nulhypothese verwerpen (een onschuldige veroordelen).
\item Type II: falen om een valse nulhypothese te verwerpen (een schuldige vrijspreken).
 \end{itemize}

\vfill\null % SOURCe: https://tex.stackexchange.com/questions/8683/how-do-i-force-a-column-break-in-a-multi-column-page
\columnbreak

\afbeelding[8cm]{img/Lrechtspraakstatistiek2.jpg}{Parallel type-I- en type-II-fout}{fig:fouten} 

Zoals in de klasactiviteit aangegeven, wordt een type-I-fout meestal als ergste gezien en de conclusie wordt zodanig getrokken dat de kans hierop beperkt is tot een zeer klein getal. Dit is het significantieniveau $\alpha$. Qua intu\"itie kun je de leerlingen meegeven dat een kleinere waarde van $\alpha$ overeenkomt met een meer behoudsgezinde (of conservatieve) houding ten opzichte van de nulhypothese en zodoende een kleinere kans op type-I-fouten. Zonder in detail te gaan, geven we nog mee dat we de type-II-fout minder onder controle hebben en dat die onder meer afhangt van de steekproefgrootte.

De bovenstaande gelijkenissen laten zien dat er in de rechtspraak elementen bestaan die analoog zijn aan de hypotheses en die ook de twee soorten van fouten kunnen verduidelijken. Eveneens geeft dit extra inzicht in het feit dat het onmogelijk is om type-I- en type-II-fouten gelijktijdig te vermijden. Ze zijn onlosmakelijk verbonden met het beslissingsproces dat gebaseerd is op beperkte informatie.

Toch is het belangrijk om te beseffen dat het hier slechts over analogie\"en gaat en dat rechtspraak niet gezien mag worden als een geval van hypothesetoetsen. Zo is het statistisch proces gebaseerd op objectieve feiten en ligt het significantieniveau $\alpha$ vast bij aanvang van de test. Bij rechtspraak zullen de feiten door zowel aanklaging als verdediging geregeld aangebracht worden vanuit een subjectief of emotioneel kader. Ook kunnen we de juryleden of rechter moeilijk zien als objectieve beslissingsmachines die los staan van menselijke emoties en werken op een vast significantieniveau. De antwoorden die je bijvoorbeeld krijgt op puntje 4 van bovenstaande lesopdracht zijn daar een mooi voorbeeld van.  

Op Khan Academy (\url{khanacademy.org}) staan enkele video's en een oefening waarbij leerlingen moeten aangeven wat een type-I- of type-II-fout is voor een gegeven hypothesetest. Ze zijn te vinden onder het onderwerp \emph{Potential errors when performing tests}.
 
\section{Betrouwbaarheidsintervallen} \label{S_5}
Bij een hypothesetoets zoals in de voorgaande paragrafen, vertrokken we steeds van een hypothese over een populatieparameter. In het voorbeeld over moraliteitsbesef bij baby's was dit de kans $p$ om voor de helper kiezen. Daarbij is het testen van de waarde $p=0,5$ een natuurlijke keuze omdat deze waarde aangeeft dat de baby's geen voorkeur hebben voor de ene of de andere pop. Bij het toetsen van hypotheses gaan we dan na hoe waarschijnlijk de geobserveerde data zijn als de nulhypothese waar is. We vragen ons met andere woorden af hoe compatibel de data zijn met de geobserveerde waarde van de populatieparameter onder de nulhypothese. %Dergelijke hypothesetoetsen kunnen ook voor populatiegemiddelden gedaan worden door gebruik te maken van een normale verdeling. 

Geregeld zijn er echter situaties waarin er geen a priori hypothese is over de populatieparameters zoals het populatiegemiddelde. Stel dat je bijvoorbeeld een onderzoek wilt doen naar de gemiddelde dagelijkse schermtijd voor Instagram bij 17-jarigen.  We hebben daar op voorhand vaak geen goed idee van. We nemen dan een steekproef uit de populatie en willen zo een idee krijgen over de waarde van het populatiegemiddelde $\mu$. %We hebben echter maar gegevens van \'e\'en steekproef. 

Een logische keuze is om $\mu$ te schatten op basis van het gemiddelde geobserveerd in de steekproef (genoteerd $\bar{x}$). Niettegenstaande deze schatting informatief is, ontbreekt er iets: ze houdt geen rekening met de variabiliteit van de steekproeftrekking. Inderdaad, als we een nieuwe steekproef nemen uit diezelfde populatie, verwachten we dat het steekproefgemiddelde zal wijzigen. Deze variabiliteit willen we terugzien bij het beantwoorden van de vraag 'wat is de waarde van $\mu$?'. Een betrouwbaarheidsinterval zal ons hier een gepast antwoord op geven.

\subsection*{Steekproevenverdeling en de centrale limietstelling}
We starten met enkele lesactiviteiten waarbij de leerlingen kunnen wennen aan het concept van een steekproevenverdeling die beschrijft hoe bijvoorbeeld de waarde van het steekproefgemiddelde $\overline{x}$ kan vari\"eren tussen verschillende steekproeven. 

We kiezen er bij deze lesactiviteiten bewust voor om te focussen op het steekproefgemiddelde $\overline{x}$ en niet meer op een (steekproef)proportie zoals in de eerdere paragrafen. Zo bekomen we immers sneller de centrale limietstelling die ons leidt naar de normale verdeling en relatief eenvoudige formules voor betrouwbaarheidsintervallen. De steekproevenverdeling voor een proportie zou ons leiden naar binomiale verdelingen die veel minder geschikt zijn voor het opstellen van betrouwbaarheidsintervallen. We komen hier later nog even op terug.

In de eerste lesactiviteit starten we met een klasgesprek.
\end{multicols}

\clearpage

\begin{lesactiviteit}{Steekproevenverdeling}
\begin{vraagenantwoord}

    \vraag{Stel dat je de gemiddelde duur wilt weten van alle zwangerschappen voor baby's geboren in een bepaald jaar en in ons land, hoe ga je dat doen?}
    \antwoord{Je ondervraagt alle ouders of ziekenhuizen in Belgi\"e naar de duur van de zwangerschap.}
        
    \vraag{In 2008 zijn er meer dan 128000 baby's geboren in Belgi\"e (gegevens van StatBel). Het is dus nogal veel werk om al die mensen te bevragen. Hoe los je dit op?}
    \antwoord{Je ondervraagt een kleinere groep ouders van baby's geboren in dat jaar.}
        
    \vraag{Welke symbolen gebruiken we voor het gemiddelde van de volledige groep en het gemiddelde van die kleine groep die je bevraagt?}
    \antwoord{Populatiegemiddelde $\mu$ voor de volledige groep en steekproefgemiddelde $\overline{x}$ voor de bevraagde groep.}
        
    \vraag{Zijn het populatiegemiddelde en steekproefgemiddelde gelijk?}
    \antwoord{Niet per se. Hoewel het mogelijk is dat beide dezelfde waarde hebben, is dit meestal niet het geval. Het steekproefgemiddelde varieert ook van steekproef tot steekproef.}
        
\end{vraagenantwoord}
\end{lesactiviteit}
\begin{multicols}{2}
Merk op dat we $\bar{x}$ noteren voor het gemiddelde geobserveerd in de enkele steekproef
die we getrokken hebben. De variabele die alle mogelijke waarden van het steekproefgemiddelde alsook de bijhorende kansen beschrijft, noteren we als $\overline{X}$. 

Voor het vervolg van deze les laten we de leerlingen zelf kennis maken met de variatie in $\overline{X}$. Daarbij willen we de verdeling aanbrengen vanuit een voorbeeld zoals we eerder al deden bij de binomiale verdeling.
\end{multicols}

\begin{lesactiviteit}{Steekproevenverdeling: vervolg 1}

 Trek nu een steekproef uit de database met de duur van alle zwangerschappen in Vlaanderen via \url{http://www3.uhasselt.be/dbgeboorten/default.asp}. Selecteer een steekproefgrootte van 30 en het jaartal 2008. Je selecteert de kolom van duur van de zwangerschap en verwerkt deze in Excel of GeoGebra of \ldots
 
 \begin{vraagenantwoord} 
 
    \vraag{Maak een histogram voor de duur van de zwangerschap in jouw steekproef. Welke vorm heeft dit histogram?}
    \antwoord{Scheef naar links.}
    
    \vraag{Kun je deze vorm verklaren?}
    \antwoord{Dit komt door de aanwezigheid van premature baby's. Uitschieters langs rechts zijn er niet. Een zwangerschap die te lang duurt wordt in de praktijk immers kunstmatig ingeleid.}

    \vraag{Wat zou de vorm van het histogram van de populatie zijn?}
    \antwoord{Het histogram zal eveneens scheef naar links verdeeld zijn om dezelfde reden.}
    
    \vraag{Wat is de gemiddelde duur van een zwangerschap in jouw steekproef? Vergelijk dit met het steekproefgemiddelde van enkele andere leerlingen in je klas.}
    
    \antwoord{Het valt op dat de steekproefgemiddelden verschillend kunnen zijn tussen de verschillende steekproeven.}
    \vraag{Zijn de verschillen tussen jullie steekproefgemiddelden groot?}
    \antwoord{Nee, de variatie in deze resultaten is klein, veel kleiner dan de variatie in de duur van de individuele zwangerschappen.}
    
\end{vraagenantwoord}
\end{lesactiviteit}
\begin{multicols}{2}
Het is nuttig om deze lesactiviteit te be\"eindigen door een overzicht te maken van alle bekomen gemiddelden van je leerlingen. Het valt dan op dat deze bekomen verdeling niet/veel minder scheef naar links is en meer symmetrisch. De belangrijkste conclusie is alvast dat verschillende steekproeven aanleiding kunnen geven tot een verschillend steekproefgemiddelde, ook al was de populatie en steekproefgrootte voor iedereen hetzelfde.
     
Je kunt deze lesactiveit eveneens verderzetten in een onderzoek dat nog dieper ingaat op de steekproevenverdeling en in essentie de centrale limietstelling introduceert via een website die in Uitwiskeling 34/1 al uitvoerig behandeld werd.
\end{multicols}
\begin{lesactiviteit}{Steekproevenverdeling: vervolg 2}
 Ga naar \url{https://onlinestatbook.com/stat_sim/sampling_dist/} en klik links op Begin.
	
 \begin{vraagenantwoord}
            
    \vraag{We werken verder met een populatie zoals die van de duur van een zwangerschap uit de vorige les. Welke vorm had die verdeling?}
    \antwoord{Scheef naar links.}
    
    \vraag{Teken in de eerste grafiek in de applet een verdeling die scheef naar links verdeeld is. Selecteer steekproefgrootte $n=2$ naast de derde grafiek en klik vervolgens op `animated' naast de tweede grafiek. (Let op: de applicatie gebruikt een hoofdletter $N$ voor de steekproefgrootte.) In de tweede grafiek verschijnen de steekproefwaarden, in de derde grafiek het gemiddelde hiervan. Klik nog enkele keren op ``animated''. Wat valt je op?}
    \antwoord{Verschillende steekproeven kunnen verschillende gemiddelden geven. }
    
    \vraag{We mogen dus geen blind vertrouwen hebben in het gemiddelde van \'e\'en steekproef wanneer we op zoek gaan naar het gemiddelde van de volledige populatie. Daarom gaan we als gedachtenexperiment dit proces 10000 keer herhalen. Klik daarvoor op 10000 naast de tweede grafiek om 10000 keer een steekproefgemiddelde van 2 randomwaarden te genereren. Op de derde grafiek ontstaat nu een overzicht van alle mogelijke steekproefwaarden. Vergelijk het gemiddelde van deze derde grafiek met het gemiddelde van de populatie. Zijn deze gelijk aan elkaar?}
    \antwoord{Het populatiegemiddelde is ongeveer gelijk aan het gemiddelde van de steekproefgemiddelden.}
    \vraag{Klik nu nog enkele malen op 100000 en vergelijk de gemiddelden opnieuw. Wat valt je op?}
    \antwoord{Als we heel erg veel/alle steekproeven bekijken, dan zijn deze waarden aan elkaar gelijk: $\mu =\mu_{\overline{X}}$; daarom noemen we $\overline{X}$ een onvertekende schatter van $\mu$.}

    \clearpage
    
    \vraag{Vergelijk ook de standaardafwijkingen van de eerste en derde grafiek. Zijn deze gelijk aan elkaar?}
    \antwoord{Neen, de derde grafiek heeft een kleinere standaardafwijking. ($\sigma_{\overline{X}}<\sigma$)}
    
    \vraag{Herhaal dit proces en beantwoord de bovenstaande vragen voor steekproefgrootte $n=16$ en $n=25$. Beantwoord de bovenstaande vragen over de gemiddelden opnieuw.}
    \antwoord{$\mu = \mu_{\overline{X}}$ blijft gelden. De standaardafwijking wordt kleiner naarmate $n$ groter wordt.}
                
    \vraag{Bereken de verhouding tussen de populatiestandaardafwijking $\sigma$ en de standaardafwijking van de steekproefgemiddelden $\sigma_{\overline{X}}$. Zie je een verband met de steekproefgrootte $n$?}
    \antwoord{Bij steekproefgrootte $n=16$ vinden we een verhouding van $4:1$, bij $n=25$ zien we een verhouding van $5:1$. We concluderen $\sigma_{\overline{X}} = \dfrac{\sigma}{\sqrt{n}}$.}
    
    \vraag{Hoe varieert de vorm van de derde grafiek naarmate dat $n$ groter wordt?}
    \antwoord{Naarmate dat $n$ groter wordt, wordt deze grafiek alsmaar meer symmetrisch en klokvormig. (Omdat de klokvorm niet altijd goed te herkennen is door een beperkt aantal balken in het histogram biedt de website een optie `fit normal'.)}
    
\end{vraagenantwoord}
\end{lesactiviteit}

\begin{multicols}{2}

De conclusie van de bovenstaande lesactiviteit legt een link naar de centrale limietstelling
\begin{kader}{}{Centrale limietstelling (zonder details)}
Voor grote waarden van $n$ geldt
\[ \overline{X} \sim N \left( \mu, \frac{\sigma}{\sqrt{n}}\right)\] 
\end{kader}
waarbij $\sim$ staat voor 'heeft als verdeling'.
Om een nog beter gevoel te krijgen voor de centrale limietstelling kun je de bovenstaande lesactiviteit ook laten herhalen voor andere vormen van populatieverdeling. Het is dan aangeraden om ook eens te starten met een normaal verdeelde populatie om in te zien dat $n$ dan niet per se groot moet zijn. Het bekijken van meerdere populatieverdelingen kan (afhankelijk van de snelheid en grootte van je klasgroep) wel een of meerdere extra lessen vragen. 

De eindconclusie is dat de normale verdeling van $\overline{X}$ gebruikt kan worden bij voldoende grote steekproeven (in de praktijk steekproefgroottes van minstens 30) of bij steekproeven bekomen uit normaal verdeelde populaties. 





\subsection*{Betrouwbaarheidsintervallen voor $\mu$ opstellen met de hand}

De leerlingen hebben op dit moment een zekere basiskennis opgebouwd over steekproevenverdelingen. Ze weten dat een steekproefgemiddelde kan afwijken van het populatiegemiddelde en van de gemiddelden van andere steekproeven. De leerlingen weten ook dat er in deze chaos nog enige regelmaat te vinden is. Zo geeft de centrale limietstelling inzicht in de mogelijke waarden die een steekproefgemiddelde kan aannemen.

We willen nu de kennis inzetten die de leerlingen hebben over normale verdelingen om tot betrouwbaarheidsintervallen te komen. Het is dan ook belangrijk dat de leerlingen eerst wiskundige eindterm 6.13 hebben behandeld alvorens statistische eindterm 6.1.2 behandeld kan worden. Voor scholen waar wiskunde en statistiek als twee afzonderlijke vakken gegeven zouden worden, dienen hier dus afspraken gemaakt te worden.

\clearpage

De zogenaamde 68-95-99,7\%-regel voor normale verdelingen zal een cruciale rol spelen:
	\begin{kader}{}{68-95-99,7\% regel}
            Als $X \sim N(\mu,\sigma)$, dan geldt:
            \begin{itemize}
                \item $\textrm{P}(\mu-\sigma \leq X \leq \mu+\sigma) \approx 0,68$
                \item $\textrm{P}(\mu-2\sigma \leq X \leq \mu+2\sigma) \approx 0,95$
                \item $\textrm{P}(\mu-3\sigma \leq X \leq \mu+3\sigma) \approx 0,997$
            \end{itemize}
        \end{kader}
        
We willen deze regel toepassen op de normaal verdeelde variabele $\overline{X}$ om een zinvolle schatting te kunnen maken voor een onbekend populatiegemiddelde $\mu$.
\columnbreak
Rechtstreekse toepassing van de regel geeft

\[ \textrm{P}\left(\mu-2\frac{\sigma}{\sqrt{n}} \leq \overline{X} \leq \mu+2\frac{\sigma}{\sqrt{n}}\right) \approx 0,95. \]

Bij betrouwbaarheidsintervallen is het echter de gewoonte om niet met ongeveer 95\% maar met exact 95\% te werken. De $z$-scores van de normale verdeling laten ons toe de formule aan te passen tot
\[ \textrm{P}\left(\mu-1,96\frac{\sigma}{\sqrt{n}} \leq \overline{X} \leq \mu+1,96\frac{\sigma}{\sqrt{n}}\right) = 0,95. \] De $z$-score van 1,96 zou al besproken moeten zijn bij de wiskundelessen over de normale verdeling. Die kun je terugvinden in GeoGebra, door eigenlijk 2,5\% van beide staarten af te knippen.
\end{multicols}

\afbeelding{img/LGeogebrazscore2.jpg}{Berekenen van z-scores via GeoGebra.}{fig:zscoreGeogebra}

\begin{multicols}{2}
Voor 95\% van de steekproeven geldt dus dat $\overline{x}$ tussen $\mu - 1,96\dfrac{\sigma}{\sqrt{n}}$ en $\mu + 1,96\dfrac{\sigma}{\sqrt{n}}$ ligt. Met andere woorden: voor 95\% van de steekproeven is het verschil tussen het steekproefgemiddelde $\overline{x}$ en het (onbekende) populatiegemiddelde $\mu$ niet meer dan $1,96 \cdot \dfrac{\sigma}{\sqrt{n}}$. Daaruit volgt het 95\% betrouwbaarheidsinterval voor een (onbekende) populatiewaarde $\mu$ in het kader hiernaast.
\begin{kader}{Formule}{95\% betrouwbaarheidsinterval voor $\mu$ bij gekende $\sigma$}
\begin{equation}
\left[ \overline{x} - 1,96 \frac{\sigma}{\sqrt{n}}, \overline{x} + 1,96 \frac{\sigma}{\sqrt{n}} \right] 
\end{equation}
\end{kader}

We geven verderop enkele vraagjes over de interpretatie van een betrouwbaarheidsinterval.

\end{multicols}

\clearpage

\afbeelding[14cm]{img/UWhagedissen.jpg}{}{}

\begin{lesactiviteit}{Interpretatie van een betrouwbaarheidsinterval}
 De formule van een betrouwbaarheidsinterval hangt van verschillende symbolen/getallen af. We gaan in op de interpretatie ervan.
 
 \begin{vraagenantwoord} 
    \vraag{Waaraan kan je zien dat het steekproefgemiddelde $\overline{x}$ het startpunt is voor het opstellen van het interval?}
    \antwoord{De formules zijn van de vorm $\overline{x}+\ldots$ en $\overline{x}-\ldots$}
    \vraag{De foutenmarge hangt af van verschillende elementen. Welke steekproeven geven de smalste intervallen? Grote of kleine steekproeven?}
    \antwoord{Hoe groter de steekproefgrootte $n$, hoe kleiner de foutenmarge, want $n$ staat in de noemer.}
    \vraag{Welk getal moet je aanpassen als je met een 99\% betrouwbaarheidsniveau wenst te werken in plaats van 95\%?}
    \antwoord{De $z$-waarde 1,96 moet aangepast worden (naar 2,58).}
    \vraag{De populatiestandaardafwijking $\sigma$ komt ook terug in de formule. Welk effect heeft die op het interval?}
    \antwoord{Hoe groter $\sigma$, hoe breder het interval en dus hoe groter de foutenmarge. Dit is logisch want meer variabiliteit in de populatie zorgt voor onzekerdere voorspellingen van het populatiegemiddelde. Eveneens kun je daaraan zien dat deze formule enkel bruikbaar is wanneer de populatiestandaardafwijking $\sigma$ gekend is. } 
\end{vraagenantwoord}
\end{lesactiviteit}
\begin{multicols}{2}

Het tweede puntje van deze lesactiviteit kan ook een rol spelen voor eindterm 6.1.2 die qua wiskunde geldt voor alle richtingen in de doorstroomfinaliteit. Deze eindterm benadrukt het belang van voldoende grote steekproeven voor het maken van nauwkeurige voorspellingen. Ook hier zien we dus weer dat de algemene eindtermen wiskunde uit de doorstroomfinaliteit nooit ver weg zijn.

 De betekenis van de 95\% betrouwbaarheid is voor veel leerlingen iets moeilijker te doorgronden. Net zoals bij de hypothesetoets heeft deze interpretatie betrekking op het toevalsproces. Indien we de studie vele malen herhalen en per studie het 95\% betrouwbaarheidsinterval berekenen (deze intervallen wijzigen van studie tot studie omdat de data wijzigen en dus ook het steekproefgemiddelde), dan zullen 95\% van die intervallen het werkelijke populatiegemiddelde omvatten. Ruwweg kunnen we stellen dat als 20 leerlingen elk een (verschillende) steekproef nemen en op basis hiervan een 95\% betrouwbaarheidsinterval opstellen, 19 intervallen het ware populatiegemiddelde omvatten en 1 niet. Die fout is dan geen gevolg van een foute werkwijze maar van onvermijdelijke steekproefvariabiliteit.

De bovenstaande formule laat zien dat we een betrouwbaarheidsinterval kunnen opstellen zodra we $\overline{x}$, $n$ en $\sigma$ kennen. Bij het voorbeeld over schermtijd op Instagram zou dat dus concreet willen zeggen dat je zo een interval kunt opstellen zodra je van een steekproef de steekproefgrootte en de gemiddelde schermtijd kent, alsook de standaardafwijking in schermtijd binnen de populatie. Jammer genoeg wringt het schoentje daar vaak. Bij realistische voorbeelden is het erg onwaarschijnlijk dat we de populatiestandaardafwijking $\sigma$ kennen wanneer we het populatiegemiddelde $\mu$ niet kennen. Daarom zullen we gebruik moeten maken van de steekproefstandaardafwijking $s$ en zal de kritische $z$-waarde 1,96 vervangen moeten worden door een zogenaamde kritische $t$-waarde die we bekomen uit een $t$-verdeling met $n-1$ vrijheidsgraden. Over dat laatste vind je meer informatie in Uitwiskeling 34/1. 

Het spreekt voor zich dat het gebruik van $t$-verdelingen een extra moeilijkheidsgraad met zich meebrengt. Dat is ook de reden waarom de meeste wiskundeboeken, zelfs voor richtingen met veel uren wiskunde, op dit moment vooral (onrealistische) oefeningen bevatten waarbij $\sigma$ gekend is. We zijn echter van mening dat het belangrijk is om ook realistische voorbeelden te beschouwen waarbij $\sigma$ onbekend is. De keuze om dit te berekenen met de hand of via ICT hangt daarbij best af van het doelpubliek en het aantal beschikbare lesuren.

Het met de hand berekenen van betrouwbaarheidsintervallen bij ongekende $\sigma$ kan relevant zijn in richtingen met de eindtermen van \emph{gevorderde wiskunde}, maar binnen \emph{statistiek} is het relevanter om bij onbekende $\sigma$ gebruik te maken van ICT.

\subsection*{Betrouwbaarheidsinterval voor $\mu$ opstellen met GeoGebra}

In GeoGebra Classic vind je bij Waarschijnlijkheidsrekening een tabblad dat zich volledig toelegt op betrouwbaarheidsintervallen en hypothesetesten. 

We keren daarvoor even terug naar de lesactiviteit Steekproevenverdelingen: vervolg 1. In die lesactiviteit vond een bepaalde leerling dat in een steekproef van 30 geboorten de gemiddelde zwangerschap 38,3 weken duurde. Gezien er geen gegevens zijn over de standaardafwijking van de populatie, gebruiken we de standaardafwijking van de steekproef, in dit geval 
\[s=\sqrt{\frac{\Sigma (x_i-\overline{x})^2}{n-1}}=2,6.\] 
Zo zie je verderop een berekening met GeoGebra van een 95\% betrouwbaarheidsinterval voor deze gegevens door middel van de t-verdeling. Het interval loopt van 37,3 tot 39,3 weken. 
\end{multicols}
\begin{center}
\afbeelding[17cm]{img/LGeogebrainterval3.jpg}{Intervallen opstellen via Geogebra}{fig:Geogebrainterval}
\end{center}
\begin{multicols}{2}

Als de leerlingen voornamelijk gebruik maken van ICT voor het opstellen van betrouwbaarheidsintervallen, dan dient er wel voldoende aandacht te zijn voor het nagaan van de nodige voorwaarden en voor het trekken van de juiste conclusies. Voor het al dan niet kennen van de populatiestandaardafwijking en het bijbehorende verschil tussen $z$-waarden en $t$-waarden voor het opstellen van een betrouwbaarheidsinterval schiet GeoGebra gelukkig ter hulp. Er wordt immers expliciet gevraagd naar $\sigma$ bij een $z$-schatting en naar $s$ voor een $t$-schatting.

\subsection*{Schatten van proporties}



%Bovenstaande voorbeelden over het opwerpen van 16 muntstukken of het bestuderen van werkloosheidscijfers tonen echter dat deze aanpak niet optimaal is. Zo is de populatiekans voor munt te gooien met een muntstuk perfect 50\%, maar het is goed mogelijk om in een steekproef met 16 worpen 11 keer munt te bekomen. Toch toonde datzelfde voorbeeld dat de steekproefwaarden niet eender wat konden zijn. Zo is er 95,1 \% kans dat in de steekproef met 16 munten minstens 4 en hoogstens 11 keer munt werd gegooid. Er is dus wel een afwijking van de te verwachten 8 keer munt, maar deze blijft beperkt.

%\end{multicols}
%\afbeelding[17cm]{StaafdiagrammuntGeogebra3.jpg}{95\% zekerheid in een staafdiagram.}{fig:StaafdiagramKopBinomiaalBerekening} 
%\begin{multicols}{2}


Tot slot van deze paragraaf geven we mee dat een compleet analoge redenering opgebouwd kan worden om na te gaan onder welke voorwaarden de schatter $\hat{p}$ normaal verdeeld is zodat we ook betrouwbaarheidsintervallen voor een populatieproportie $p$ kunnen opstellen. 

Het idee is dat we vertrekken van de binomiale verdeling die achter $\hat{P}$ schuilt. Met de notatie zoals in paragraaf \ref{S_4} geldt immers $\hat{P} = \displaystyle \frac1n \sum_{i=1}^n Y_i$ zodat $\hat{P} = \dfrac{T}{n}$ waarbij $T$ de binomiaal verdeelde toetsingsgrootheid uit paragraaf \ref{S_4} is. Aangezien dit echter een discrete verdeling is in plaats van een continue is het hier nagenoeg onmogelijk om een interval van waarden te bepalen waarin exact 95\% valt. Daarom worden betrouwbaarheidsintervallen voor de populatieproportie $p$ vrijwel alleen bekeken in gevallen waarbij de binomiale verdeling benaderd kan worden door een normale verdeling, zoals aangehaald in \autoref{fig:standaard2} uit paragraaf \ref{S_3}. 

Dit levert echter technische issues zoals de continu\-iteitscorrectie of het feit dat een proportie $p$ altijd ligt tussen 0 en 1, wat we bij een normale verdeling niet kunnen garanderen. Ook zijn er verschillende aanpakken voor het opstellen van de foutenmarge in dit geval. De diepgang van deze zaken lijkt ons echter aangewezen voor richtingen met meer uren wiskunde. Het enige interval dat ons haalbaar lijkt voor de richting humane wetenschappen is het zogenaamde Wald-betrouwbaarheidsinterval dat geen enkele correctie doorvoert in functie van de drie bovenstaande bedenkingen. 

    \begin{kader}{}{95\% Wald-betrouwbaarheidsinterval voor een populatieproportie $p$}
           \[  \left[ \hat{p} - 1,96 \sqrt{\frac{\hat{p}\left( 1- \hat{p}\right)}{n}}, \hat{p} + 1,96 \sqrt{\frac{\hat{p}\left( 1- \hat{p}\right)}{n}}\right] \]
    \end{kader}

Omdat in het hoger onderwijs de focus iets vaker ligt op het juist kunnen toepassen van de formules dan op het van buiten leren van de formules, lijkt het ons hier het voornaamste om aandacht te spenderen aan de nodige voorwaarden die gelden voor de benadering van een binomiale verdeling door een normale (en het daaruitvolgende Wald-betrouwbaarheidsinterval). Nuttige oefeningen zijn te vinden op Khan Academy (\url{khanacademy.org}) onder de lessenreeks \emph{Confidence intervals for proportions}. 

	
\section{Vakoverschrijdende projecten} \label{S_6}



\subsection*{Het GAISE-rapport}
    
    Er zijn verschillende manieren om statistiek te onderwijzen. Klassieke handboeken delen hun inhoud op volgens drie thema's: beschrijvende statistiek, kansrekenen en inductieve (of verklarende) statistiek. Het eerste deel wordt vaak als minst moeilijk ervaren, terwijl het tweede en derde deel uitdagender zijn. Het deel rond kansrekening levert inhouden aan die noodzakelijk zijn voor inductieve statistiek, maar dit kan er ook voor zorgen dat leerlingen de meerwaarde van statistiek niet meer begrijpen door een overdaad aan formules. Dit kan een impact hebben op de motivatie en resultaten voor het vak. Een mogelijke oplossing voor dit probleem ligt in de alternatieve volgorde van lesonderwerpen zoals we hierboven hebben voorgesteld: het is mogelijk om eerst de idee\"en van de inductieve statistiek te introduceren als motivatie voor de kansberekening.
    
    In 2016 verscheen er een rapport met richtlijnen voor het evalueren en het onderwijzen van introductiecursussen statistiek: het GAISE-rapport (Carver et al., 2016). Dit rapport wordt onderschreven door de \emph{American Statistical Association}, \'e\'en van de grootste beroepsverenigingen voor statistici wereldwijd en de Amerikaanse \emph{National Council of Teachers of Mathematics}. Het rapport is gratis beschikbaar en bevat voorbeelden van good practices en lesmateriaal in de hoop de nodige motivatie te bieden voor de leerlingen. Uiteraard gaan veel van deze idee\"en ook op voor de andere richtingen waarin eindtermen statistiek voorkomen onder de noemer uitgebreide of gevorderde wiskunde. Wij beperken ons hier tot de zes aanbevelingen: 

\begin{enumerate}
    \item Leer studenten statistisch denken, met specifieke aandacht voor: 
    \begin{enumerate}
        \item Statistiek als een probleemoplossend en besluitvormend onderzoeksproces. 
        \item Multivariabel denken. 
    \end{enumerate}
    \item Focus op conceptueel begrip en beperk kansrekening tot een minimum.
    \item Integreer echte data met context en doel. 
    \item Stimuleer activerend leren. 
    \item Maak gebruik van technologie om concepten te onderzoeken en data te analyseren. 
    \item Maak gebruik van evaluatie om het leren te bevorderen en te evalueren. Deze evaluaties moeten in lijn zijn met de leerdoelen: o.a. statistisch denken, conceptueel begrijpen, kritisch evalueren. \\
\end{enumerate}

De eerste drie aanbevelingen weerspiegelen goed wat statistiek betekent: het is de wetenschap van het leren uit data en van het meten, controleren en communiceren van onzekerheid (Davidian en Louis, 2012). In deze omschrijving staan zowel de onderzoeksvraag, het ontwerp van de studie, de data als de methoden die toelaten om te gaan met onzekerheid centraal. Deze vier elementen komen bovendien sterk overeen met de vier stappen uit een onderzoek zoals beschreven in de OVUR-methode: \textbf{O}ri\"enteren, \textbf{V}oorbereiden, \textbf{U}itvoeren en \textbf{R}eflecteren. 

Door een sterke focus te leggen op de data en hun context (de onderzoeksvragen die men wenst te beantwoorden) blijft de connectie tussen statistiek en de `echte wereld' behouden. 
Voor het doelpubliek dat beschreven wordt door de eindtermen \emph{Statistiek} is deze connectie extra relevant omdat, in tegenstelling tot de meer wiskundige richtingen, het hier niet de bedoeling is om eigenschappen te bewijzen of te experimenteren met nieuwe formules. Het einddoel is om de leerlingen inzichten te geven waardoor ze in staat zijn om statistiek te gebruiken in een realistisch onderzoek in hun verdere studie of carri\`ere. Omdat dit onderzoek veelal zal gebeuren vanuit een sociaal, cultureel of ethisch vraagstuk, kan het bijzonder nuttig zijn om vakoverschrijdende projecten op te starten waarbij statistiek instaat voor het verwerken en interpreteren van de data en een of meerdere andere vakken instaan voor het formuleren van de onderzoeksvraag.

\subsection*{Een voorbeeld uit de praktijk}
We geven hieronder een voorbeeld van hoe zo een vakoverschrijdend project vorm heeft gekregen bij een van de loepschrijvers. Dit gebeurde binnen het vak statistiek, dat in afwachting van de vernieuwde leerplannen in de complementaire uren georganiseerd wordt en de uitbreidingsdoelstellingen van het vak wiskunde volgt. Deze projecten vonden allemaal plaats in schooljaar '19-'20, '20-21, '21-'22 en '22-'23 waardoor ze spijtig genoeg gehinderd werden door het coronavirus. De leerlingen die dit schooljaar in het zesde middelbaar zitten, behoren tot de eerste lichting die het project in zijn meest uitgebreide vorm zullen kunnen afronden. 

We volgen in de beschrijving de stappen van de OVUR-methode en geven bij elke stap weer welke rol statistiek zou kunnen spelen en welke eindtermen (van andere vakken) relevant kunnen zijn. Qua algemene eindtermen denken we uiteraard aan eindterm 6.1.4 van statistiek alsook de generieke doorstroomcompetentie over onderzoeksprocessen.

        
    \begin{kader}{Eindterm}{Generieke doorstroomcompetenties}
             1.1.3 De leerlingen doorlopen een onderzoeksproces in functie van een onderzoeksvraag.
    \end{kader}

\subsection*{Ori\"enteren}
Het voorbeeldproject dat we hier als rode draad zullen gebruiken, ging over sociale media. Daarbij werden de leerlingen getriggerd via de documentaire `The social dilemma' op Netflix. Gebaseerd hierop gingen de leerlingen op zoek naar een mogelijk onderzoeksonderwerp. Hierbij suggereerden ze zelf: de invloed van socialemediagebruik op welbevinden, de invloed van socialemediagebruik op slaap, het verschil in socialemediagebruik tussen jongens en meisjes, de invloed van sociale media op aankoopgedrag \ldots We komen later terug op de formulering van deze vragen maar dit ori\"enterende gedeelte kan afhankelijk van het onderwerp in een van de andere vakken gebeuren. Voor dit voorbeeld denken we aan gedragswetenschappen of informatica.

    \begin{kader}{Eindterm}{Specifieke eindtermen sociale wetenschappen}
             15.1.3 De leerlingen reflecteren over maatschappelijke vraagstukken aan de hand van aangereikte kwalitatieve en kwantitatieve onderzoeksresultaten, met aandacht voor ruimtelijke en historische achtergronden.
    \end{kader}
   \begin{kader}{Eindterm}{Digitale competenties voor de doorstroomfinaliteit}
             4.7 De leerlingen beoordelen kritisch wederzijdse invloeden tussen enerzijds het individu en de samenleving en anderzijds media, digitale infrastructuur en digitale toepassingen.
    \end{kader}  









    %Het ori\"enterende gedeelte van het onderzoekswerk is sterk afhankelijk van de vakken die samen met statistiek in het project betrokken worden. Binnen de richting humane wetenschappen bieden de vakken sociale wetenschappen en gedragswetenschappen mooie opportuniteiten met onderstaande eindtermen.


    
    %Maatschappelijke vraagstukken kunnen hierbij over allerlei onderwerpen gaan, al zijn binnen de richting humane wetenschappen armoede, ongelijkheid, migratie, eenzaamheid en psychisch welzijn duidelijke kandidaten.            
	
    %\begin{kader}{Eindterm}{Specifieke eindtermen gedragswetenschappen}
     %        14.1.2 De leerlingen reflecteren over sociaal gedrag aan de hand van sociaalpsychologische theorieën.
    %\end{kader} 
    
    %Mogelijke onderwerpen zijn hier pestgedrag of de verschillen in gedrag thuis en op school of verschillen tussen online gedrag en offline gedrag. Dat laatste past dan weer mooi binnen de digitale competenties:

    %Vanuit het vak filosofie zouden dan weer ethische kwesties getoetst kunnen worden bij een doelpubliek.
    
    %\begin{kader}{Eindterm}{Specifieke eindtermen filosofie}
     %       5.1.2 De leerlingen reflecteren over ethische stromingen en vraagstukken aan de hand van filosofische begrippen.
    %\end{kader}

    %Zo kunnen de leerlingen nagaan welke ethische visies aanwezig zijn binnen hun leeftijdscategorie of school. Misschien denkt hun leeftijdscategorie niet hetzelfde over ethische kwesties als de gemiddelde Belg? Eveneens kunnen binnen het vak filosofie of geschiedenis politieke kwesties aan bod komen die aanleiding geven tot interessante onderwerpen om te onderzoeken.

   % Zodra de leerlingen enigszins een beeld hebben van de richting die ze willen inslaan met hun onderzoek, gaan ze over op een literatuurstudie. Welke resultaten zijn bijvoorbeeld al gekend voor de vraag die ze in gedachten hebben? In deze stap is het cruciaal dat er kritisch wordt omgegaan met de bronnen.

            
\subsection*{Voorbereiden}

Dit onderzoek over sociale media moest binnen het vak statistiek verder gespecificeerd worden tot twee onderzoeksvragen waarbij de leerkracht statistiek erop toezag dat ze haalbaar en duidelijk afgelijnd waren. Zo kan statistiek wel ingezet worden om een verband tussen twee variabelen te ontdekken, maar is het moeilijk in deze context om tot besluit van een oorzakelijk verband te komen. Daarom moest voor veel onderzoeksvragen zoals gesuggereerd door de leerlingen het woord `invloed' vervangen worden door `verband'. 

Ook voor de haalbaarheid van de analyse moet opgelet worden. Hebben de leerlingen voldoende voorkennis om deze vraag statistisch aan te pakken? Zo is de vraag `Is er een verband tussen de hoeveelheid schermtijd op je smartphone en je schoolresultaten?' wel interessant, maar is die moeilijk te onderzoeken voor leerlingen zonder voorkennis over lineaire regressie. De leerlingen die in de huidige tweede graad zitten, zouden hier volgens de vernieuwde eindtermen wel ervaring mee moeten hebben, maar je spreekt best eens met de wiskundecollega's van de tweede graad om een goed zicht te hebben op de beschikbare voorkennis.

Naast de haalbaarheid op vlak van beschikbare kennis statistiek is ook de beschikbaarheid van data relevant. Er zijn tegenwoordig veel betrouwbare websites die veel informatie bieden. We denken bijvoorbeeld aan Statbel, maar soms is het nodig om zelf data te verzamelen. Wanneer leerlingen deze data verzamelen via een enqu\^ete zijn er weer drie nieuwe aandachtspunten om te behandelen in het vak statistiek.

Het eerste aandachtspunt is het doelpubliek: een enqu\^ete die zich richt op leerlingen van de eigen school is makkelijk te organiseren, een enqu\^ete in een andere school of een andere leeftijdscategorie is vaak moeilijker. Durf hier te spreken uit jouw ervaring. Het heeft weinig zin om ambitieuze doelen te stellen waarvan je op voorhand weet dat ze niet haalbaar zijn. 

Het tweede aandachtspunt is de selectie van de steekproef. Deze kan volledig aselect bekomen worden of gestratificeerd. Een aselecte steekproef heeft als voordeel dat dit mooi in lijn ligt met de steekproevenverdeling zoals besproken in de vorige paragrafen. Er zijn echter enkele nadelen: zo is het niet altijd evident om één groot bestand te hebben waarin alle leerlingen van de school opgelijst worden. Het gebrek van zo een lijst kan het moeilijk maken om te beseffen dat `leerling 112' eigenlijk overeenkomt met leerling 4 uit klas 2B. Een gestratificeerde steekproef waarbij we uit elke klas een vast percentage van de leerlingen selecteren is praktisch een pak haalbaarder omdat je dan enkel klaslijsten nodig hebt. Ook geeft dit het voordeel dat het selectiewerk verdeeld kan worden over verschillende leerlingen. Deze aanpak heeft echter het technische nadeel dat sommige veronderstellingen zoals onafhankelijkheid of \'e\'en verdeling voor de volledige populatie in het gedrang komen. Ongeacht welke keuze er wordt gemaakt, dient de representativiteit bewaakt te worden.
            
    \begin{kader}{Eindterm}{Eindtermen wiskunde voor de doorstroomfinaliteit}
            6.12 De leerlingen verklaren het belang van randomisatie en representativiteit bij steekproeven en bij experimenteel vergelijkend onderzoek voor het formuleren van statistische besluiten over een populatie.
    \end{kader}

Het derde aandachtspunt is de duidelijkheid van de enqu\^etevragen en de bijhorende antwoordmogelijkheden. Je zult merken dat leerlingen hier vaak te kort door de bocht gaan en dat bijsturing nodig is. Zo zijn veel ja/nee-vragen vaak beter te stellen op een schaal van helemaal akkoord tot helemaal niet akkoord of van altijd tot nooit om de nodige nuance te leggen. Ook wijs je de leerlingen op het gevaar van suggestieve vragen.

In functie van de ICT-doelstellingen is het hier bovendien handig als jij (of indien van toepassing: de leerkracht ICT) ingaat op de mogelijkheden van online bevragingen via bijvoorbeeld Google Forms. De laatste stap in de voorbereidende fase ligt waarschijnlijk bij de leerkracht. Je zult bijvoorbeeld alle enqu\^etevragen voor een bepaalde populatie bundelen en essentieel identieke vragen zo goed mogelijk samenvoegen. Anders riskeren we immers dat respondenten afhaken bij de enqu\^etevragen die later komen in de lange reeks. 

Een van de geselecteerde vragen bij het praktijkvoorbeeld over sociale media was de volgende: `Is er een verband tussen geslacht en inschattingsvermogen van schermtijd?' Om deze vraag in een enqu\^ete te plaatsen, was  verduidelijking over het begrip schermtijd nodig. Willen we het hebben over schermtijd op sociale media of in het algemeen? Wanneer we de schermtijd op sociale media onderzoeken, dan moet het voor iedereen duidelijk zijn wat daarmee bedoeld wordt. Behoort YouTube bijvoorbeeld tot de sociale media? Daarom is het voor een enqu\^ete duidelijker om specifiek de schermtijd via smartphone voor \'e\'en bepaalde app te vragen of simpelweg de totale schermtijd voor alle apps tesamen. In dit praktijkvoorbeeld werd gekozen voor de tweede optie. 
    
Daarnaast was het voor deze vraag cruciaal dat de bevraagden wisten waar ze deze schermtijd konden terugvinden en in welke eenheid (minuten per dag/ uren per week/ \ldots) geantwoord moest worden (we kozen voor uren per week). Er werd dan ook de nodige uitleg voorzien in de enqu\^ete. Omdat er ondanks de uitleg personen konden zijn die hun schermtijd niet kunnen terugvinden, moesten we de optie toevoegen \emph{Ik heb deze schermtijd niet teruggevonden}. Deze personen zouden later als non-response ge\"interpreteerd worden. 
    
Voor de bevraging van het geslacht werd gekozen voor de volgende formulering: Met welk geslacht identificeer je je? Daarbij werden de opties M/V/X gegeven.
    
Naast dit voorbereidende werk in het statistiekvak kan er eveneens in andere vakken aandacht besteed worden aan de literatuurstudie. Dat kan plaatsvinden in het vak dat het onderzoeksthema voorgedragen heeft, maar ook in het vak geschiedenis waar reeds aandacht besteed moet worden aan de betrouwbaarheid van bronnen.
    
\begin{kader}{Eindterm}{Generieke doorstroomcompetenties}
             1.1.2 De leerlingen beoordelen informatie kritisch op bruikbaarheid, validiteit en betrouwbaarheid in functie van hun informatiebehoefte.
\end{kader}
%\begin{kader}{Eindterm}{Historisch bewustzijn binnen de doorstroomfinaliteit}
        %     8.4 De leerlingen evalueren de presentatie, de context, de betrouwbaarheid, de representativiteit en de bruikbaarheid van historische bronnen in het licht van een historische vraag.
%\end{kader}

Het feit dat er binnen het vak geschiedenis meer gefocust wordt op historische bronnen hoeft hierbij geen belemmering te vormen. Veel algemene principes voor het detecteren van betrouwbare bronnen zijn ook in een hedendaagse context nuttig. Ook zal de leerkracht geschiedenis vanuit diens opleiding vaak de expert ter zake zijn betreffende bronnenmateriaal. Tot slot geven we mee dat de leerlingen, in functie van een later verslag, hun bevindingen uit de literatuurstudie best nu al ergens neerschrijven.



    %Vermoedelijk is er nog geen onderzoek gedaan naar ethische kwesties over DNA of zelfrijdende auto's of sociale media binnen jouw school. Belangrijk bij het opstellen van een onderzoeksvraag is na te denken over welke populatie we uitspraken wensen te doen. Vaak zullen we binnen de sociale wetenschappen gebruik maken van een "gemakssteekproef", en dienen we ons de vraag te stellen voor welke populatie deze representatief is. 


    
    %Gezien deze eindterm geformuleerd staat binnen het vak wiskunde, kan deze stap van het werk door de betreffende wiskundecollega georganiseerd worden. Een mogelijke werkwijze om deze steekproef te kiezen is om te gaan voor een gestratificeerde steekproef waarbij eerst wordt bepaald hoeveel leerlingen per klas/richting/jaar/gender/\ldots geselecteerd moeten worden. De strata of subpopulaties worden daarbij best slim gekozen ten opzichte van de onderzoeksvragen. Voor deze keuze sta jij in als leerkracht, al baseer je je hier op de voorstellen van de leerlingen. Tenzij er redenen zijn om anders te kiezen, raden we aan om de klassen te zien als de strata. Zo vermijd je te grote getallen in de random number generator uit onderstaande lesactiviteit. Ook vermijd je zo het risico dat de verschillende strata verschillende verdelingen kennen binnen de populatie, hetgeen een probleem zou zijn gezien de achterliggende theorie slechts \'e\'en verdeling veronderstelt voor de volledige populatie. De te gebruiken steekproefgrootte wordt eveneens best bepaald door de leerkracht, zowel met het oog op representativiteit als op de voorwaarden die nodig zijn voor latere berekeningen.
            
    %Voor onderstaande lesactiviteit groepeer je eerst de leerlingen per doelpubliek/populatie. Deze leerlingen zullen nu samen een gestratificeerde steekproef cre\"eren van de populatie. De leerlingen die hun data online kunnen vinden, verdeel je over de groepen zodat ook zij meewerken bij deze activiteit. De leerlingen krijgen dan per twee of drie enkele strata toegewezen in de vorm van anonieme klaslijsten waarbij enkel noodzakelijke informatie zoals jaar, studierichting, gender, \ldots vermeld is.
%\end{multicols}

%\begin{lesactiviteit}{Cre\"eren van een gestratificeerde steekproef}
        
	%\begin{vraagenantwoord}
         %   \vraag{Er zijn $1200$ personen in deze school (= de populatie), we zullen een steekproef nemen van grootte $n=100$. Als we gestratificeerd werken, hoeveel personen moeten we dan selecteren uit een klas (= een stratum) van $25$ personen? (Vervang in deze vraag 1200, 100 en 25 door de grootte van de populatie, steekproef en stratum in jouw specifieke situatie. Werk eventueel algemeen met letters $N, n, x$ alvorens concrete getallen te gebruiken).}
          %  \antwoord{In elk stratum selecteren we een relatieve frequentie van $\dfrac{100}{1200}=\dfrac{1}{12}$. Voor een stratum van $25$ personen geeft dat dus $\dfrac{25}{12} \approx 2,08$ personen. In de praktijk selecteren we dus 2 personen.}
           % \vraag{Gebruik nu een random number generator om de volgnummers van 2 personen te selecteren. Geef de lijst met geselecteerde volgnummers aan de leerkracht.}
         %   \antwoord{Google heeft een ingebouwde random number generator. Je bekomt deze door \emph{`random number generator'} in te geven bij de zoekbalk. Je kunt ook \url{random.org} gebruiken. Die website kun je gebruiken voor simulaties van random nummers, dobbelsteenworpen, muntworpen, speelkaarten, \ldots}
       % \end{vraagenantwoord}

%\end{lesactiviteit}
%\begin{multicols}{2}
        
    

          
\subsection*{Uitvoeren}
    Na al het voorbereidende werk is het tijd voor de uitvoering. De data worden bekomen via enqu\^etes bij de steekproeven of via online bronnen. Bij de enqu\^etes met vragen van meerdere leerlingen kun je ervoor kiezen om alle antwoorden met de leerlingen te delen of enkel de antwoorden op de vragen die voor hen relevant zijn. Beide aanpakken kennen hun voor- en nadelen. Door de leerlingen alle antwoorden te geven, worden ze getraind in het filteren van informatie via bijvoorbeeld de filterfunctie in Excel. Anderzijds zou je omwille van privacyredenen kunnen kiezen om enkel de relevante antwoorden te delen met de leerlingen. 

    Nu volgt het moment waarop de theorie van statistiek in werking treedt. Bij de voorbeeldvraag `Is er een verband tussen geslacht en inschattingsvermogen van schermtijd?' vergeleken de leerlingen daarbij de personen die zich als man identificeerden met deze die zich als vrouw identificeerden. In de enqu\^ete waren er geen non-binaire genderidentificaties, hetgeen het onderzoek iets makkelijker maakte gezien er nu maar twee groepen vergeleken moesten worden. De variabele die onderzocht werd was de verschilvariabele: \emph{Effectieve schermtijd min geschatte schermtijd}. We noteren $\mu_M$ en $\mu_V$ voor de populatiegemiddelden voor deze verschilvariabele voor beide genders. Op basis van de literatuurstudie en het algemene vermoeden van de leerlingen dat mannen nonchalanter zijn en vaker negatieve effecten (zoals het tijdverslies door schermtijd op sociale media) minimaliseren dan vrouwen, stelden de leerlingen de volgende hypotheses op:
    
\[
H_0: \mu_M = \mu_V \quad \text{versus} \quad H_A: \mu_M > \mu_V. 
\]

    Hoewel de leerlingen enkel ervaring hadden met het onderzoeken van hypotheses over \'e\'en gemiddelde en niet over het verschil van gemiddelden, kon hier toch het een en ander opgelost worden door hun kennis over P-waarden en ICT. Op GeoGebra bestaat er (in hetzelfde tabblad als getoond in \autoref{fig:Geogebrainterval}) immers de mogelijkheid om P-waarden uit te rekenen voor een gegeven test. Hier vonden de leerlingen twee mogelijkheden: een $t$-test en een $z$-test voor het verschil van gemiddelden. Het verschil tussen een $z$-test en $t$-test kenden ze reeds uit hypothesetoetsen voor het gemiddelde van een normaal verdeelde populatievariabele.
    
    Aangezien we hier niet mogen veronderstellen dat de populatiestandaardafwijkingen gekend zijn, moeten we een $t$-test gebruiken en de steekproefwaarden gebruiken, naar analogie met de betrouwbaarheidsintervallen zoals besproken in \autoref{S_5}. 
        
    De achterliggende $t$-verdeelde toetsingsgrootheid was voor de leerlingen irrelevant en te moeilijk om te bespreken. Het was hier wel nodig om te benadrukken dat deze berekeningen enkel gedaan kunnen worden in de veronderstelling dat de populatievariabelen normaal verdeeld zijn.

    De leerlingen ontdekten enkele verschillen: in de steekproef bleek dat jongens hun schermtijd vaker onderschatten dan meisjes ($\overline{x}_M=6,11$ versus $\overline{x}_V=1,95$, uitgedrukt in uren per week) maar dit verschil was niet significant op het 5\% niveau. De $P$-waarde van 0,0539 of 5,39\% liet dus wel een vermoeden achter maar nog geen bewijs, gezien deze groter is dan de significantiewaarde $\alpha=0.05$ (zie stap 5 in paragraaf \ref{S_4}). Hierbij dienen twee zaken opgemerkt te worden. Enerzijds konden verschillende deelnemers in de enqu\^ete hun totale schermtijd niet terugvinden waardoor er best veel non-response was en de bruikbare steekproefgroottes klein waren. 

    \afbeelding[7cm]{img/LTtestPwaardeGeogebra.jpg}{P-waarden rechtstreeks vinden via GeoGebra}{fig:TtestPwaardeGeogebra} 

   Anderzijds was er veel variatie op de steekproefwaarden. Dit gaf na afloop dus zeker stof tot reflecteren: bijvoorbeeld hoe kan die non-response lager gehouden worden.
   
    Na deze berekeningen gaan de leerlingen over op het schrijven van een verslag. Dit verslag begint met de motivatie van de onderzoeksvraag, gevolgd door resultaten uit de literatuurstudie, een korte beschrijving van de enqu\^ete, resultaten van de enqu\^ete en ten slotte de interpretatie van deze resultaten in functie van de onderzoeksvraag. Bij het schrijven van dit verslag is een samenwerking met de leerkracht Nederlands een meerwaarde.

    \begin{kader}{Eindterm}{Eindtermen Nederlands binnen de doorstroomfinaliteit}
            2.8 De leerlingen vatten een geschreven tekst schriftelijk samen in functie van doelgerichte informatieverwerking en communicatie.
            
            2.9 De leerlingen produceren schriftelijke en mondelinge teksten in functie van doelgerichte communicatie.
    \end{kader}
        
    Zo wordt, in een ideale wereld, het finale verslag door drie leerkrachten nagelezen: de vakleerkracht waaruit het oorspronkelijke onderwerp kwam, de leerkracht statistiek en de leerkracht Nederlands.

    Tot slot kun je ervoor kiezen om de resultaten ook te presenteren via een powerpoint of postervoorstelling. Zo worden nog enkele digitale competenties bereikt.   Een mondelinge presentatie lukte in het jaar van het onderzoek rond sociale media jammer genoeg niet, maar in een ander project presenteerden de leerlingen hun resultaten op een A3-poster waarbij minstens \'e\'en diagram getoond moest worden. De resultaten hingen gedurende een week op een zichtbare plaats op de overdekte speelplaats aangezien de resultaten interessant konden zijn voor de volledige school. 

 
    \begin{kader}{Eindterm}{Digitale competenties voor de      doorstroomfinaliteit}
            4.3 De leerlingen gebruiken doelgericht en adequaat functionaliteiten van digitale infrastructuur en toepassingen om digitaal te communiceren, samen te werken en te participeren aan initiatieven.
    \end{kader}
        

    
\subsection*{Reflecteren}
    De fase van reflecteren kunnen we in twee delen opsplitsen.
            
    Enerzijds kunnen de leerlingen reflecteren over de bekomen uitkomsten binnen het samenwerkende vak. Zou er bijvoorbeeld een reden kunnen zijn waarom we geen significant verschil konden vaststellen tussen de jongens en meisjes in deze specifieke steekproef of zou dat slechts toeval zijn? Het samenwerkende vak kan zich daarbij vooral richten op mogelijke verklaringen.
    
    %Stel bijvoorbeeld dat de ethische visie bekomen uit een steekproef op school behoorlijk verschilt van de voorspelde visies uit eerdere onderzoeken: waaraan zou dit dat kunnen liggen? Zou dit een gevolg zijn van de leeftijd of de etniciteit van de lokale bevolking?
    Binnen de lessen statistiek kan dan weer gereflecteerd worden over de werkwijze. Zeker bij onverwachte resultaten kan dit nuttig zijn. Was de steekproef bijvoorbeeld wel representatief genoeg? Zijn we zeker dat alle nodige voorwaarden voor het gebruiken van een normale of binomiale verdeling voldaan waren? Zou een andere werkwijze beter geschikt kunnen zijn voor dit onderzoek? Als leerkracht geeft dit laatste je ook de kans om de rol van statistiek in het hoger onderwijs alsook de verschillende onderwerpen die erin aan bod komen te belichten in de klas.
    
\subsection*{Niet alleen vakoverschrijdend, maar ook jaaroverschrijdend?}


    De ervaring leert dat als je het project echt vakoverschrijdend wilt aanpakken, dit best ook jaaroverschrijdend gebeurt waarbij de OVUR-methode splitst in OV in het vijfde jaar en UR in het zesde jaar. De huidige zesdejaarsleerlingen begonnen bijvoorbeeld aan het project in januari 2022 en presenteren hun werk in april 2023. Ze voerden hun literatuurstudie uit en kregen bij de lessen geschiedenis uitleg over bronnen en binnen Nederlands over bronvermelding (volgens de APA-normen die de standaard vormen binnen de sociale wetenschappen). Hun enqu\^etevragen werden in het eerste semester van dit schooljaar opgesteld en de enqu\^eteresultaten werden in januari en februari verwerkt. Op dit moment is het dus al uitkijken naar de nakende presentaties waarop deze leerlingen hun bevindingen uit de doeken zullen doen.

    %Achter het volgende %-teken vinden jullie een ander voorbeeld dat hier eerst stond maar waarbij lineaire regressie een belangrijkere rol speelt
            
    %De leerlingen hadden toen al ervaring met lineaire regressie waardoor een van de onderzoeksvragen als volgt luidde ``Heeft het gebruik van sociale media een invloed op je slaap?''. In de voorbereidende fase werd deze vraag aangescherpt. Wanneer we de schermtijd op sociale media onderzoeken, dan moet het voor iedereen duidelijk zijn wat daarmee bedoeld wordt. Behoort YouTube bijvoorbeeld tot de sociale media? Daarom is het voor een enqu\^ete duidelijker om specifiek de schermtijd via smartphone voor 1 bepaalde app te vragen, bijvoorbeeld TikTok. Daarbij ging er ook aandacht naar het feit dat sommige ondervraagden misschien uitleg moesten krijgen over hoe ze deze schermtijd konden oproepen. Zowel voor de schermtijd als de hoeveelheid slaap werd ook nagedacht over de vraagstelling (is het verschil tussen weekdagen en werkdagen relevant? Wij besloten om een weekgemiddelde te vragen.)
    
    %De steekproef werd gestratificeerd opgesteld zoals in bovenstaande lesopdracht en de geselecteerde leerlingen vulden de enqu\^ete met deze vraag en enkele andere vragen in tijdens vastgelegde uren in de refter en via Google Forms. De statistiekleerlingen kregen om wille van privacyredenen enkel de resultaten van kolommen die relevant konden zijn. 

    

            
    %Na de nodige berekeningen verschenen in het verslag een spreidingsdiagram met berekening van correlatieco\"effici\"ent en determinatieco\"effici\"ent. Ook werden betrouwbaarheidsintervallen opgesteld voor de gemiddelde schermtijd en gemiddelde hoeveelheid slaap. 
    
  

            

\end{multicols}

\stopcontents[inhoudloep]		% om inhoud voor het loep artikel te eindigen

\clearpage

%% BRONNEN %%%%%%%%%%%%%%%%%%%%%%%%%%%%%%%%%%%%%%%%%%%%%%%%%%%%%%%%%%%%%%%%%%%%%
\begin{bronnen}

    \item Carver, R., Everson, M., Gabrosek, J., Horton, N., Lock, R., Mocko, M., ... en Wood, B. (2016). \emph{Guidelines for assessment and instruction in statistics education (GAISE) college report 2016.} \url{https://www.amstat.org/docs/default-source/amstat-documents/gaisecollege_full.pdf}
    \item Davidian, M., Louis, A.T. (2012). Why Statistics? \emph{Science, 336(6077)}, 12.
    \item Hamlin, J. K., Wynn, K., \& Bloom, P. (2007). Social evaluation by preverbal infants. \emph{Nature, 450(7169)}, 557-559.
    \item Maenhout, K. (2021), \emph{Humane wetenschappen mag niet langer kneusje aso zijn} in: \emph{De Standaard}, \url{https://www.standaard.be/cnt/dmf20210905_97611830}
    \item Rollinson, C. (2023). \emph{Coin Flip - Many Trials}. \url{https://www.geogebra.org/m/yk75zwrq#material/v5ceaza5}
    \item Rossman, A. J. (2008). Reasoning about Informal Statistical Inference: One Statistician's View. \emph{Statistics Education Research Journal, 7(2)}.
    \item Siew,  C.S.Q., McCartney, M.J., Vitevitch M.S. (2019), \emph{Using Network Science to Understand Statistics Anxiety Among College Students} in \emph{Scholarship of Teaching and Learning in Psychology 5(1)}, 75–89.
    \item StatBel  (2023). \emph{Evolutie van het aantal geboorten in Belgi\"e 1830-2021}. \url{https://statbel.fgov.be/nl/themas/bevolking/geboorten-en-vruchtbaarheid#figures}
    \item Vlaamse Regering (2019). \emph{Regeerakkoord van de Vlaamse Regering 2019-2024}, \url{https://www.vlaanderen.be/publicaties/regeerakkoord-van-de-vlaamse-regering-2019-2024}
    \item Vlaamse Regering (2022). \emph{Decreet over de instroom en het optimaliseren van de studie-efficiëntie in het hoger onderwijs en overige organisatorische aspecten van het hoger onderwijs} in \emph{Belgisch staatsblad, publicatie 2022-08-24}.
    \item Wason, P.C., Johnson-Laird, P.N. (1972).
    \emph{Psychology of reasoning: Structure and content}, Harvard University Press.
\end{bronnen}

